% Options for packages loaded elsewhere
\PassOptionsToPackage{unicode}{hyperref}
\PassOptionsToPackage{hyphens}{url}
\documentclass[
]{article}
\usepackage{xcolor}
\usepackage[margin=1in]{geometry}
\usepackage{amsmath,amssymb}
\setcounter{secnumdepth}{-\maxdimen} % remove section numbering
\usepackage{iftex}
\ifPDFTeX
  \usepackage[T1]{fontenc}
  \usepackage[utf8]{inputenc}
  \usepackage{textcomp} % provide euro and other symbols
\else % if luatex or xetex
  \usepackage{unicode-math} % this also loads fontspec
  \defaultfontfeatures{Scale=MatchLowercase}
  \defaultfontfeatures[\rmfamily]{Ligatures=TeX,Scale=1}
\fi
\usepackage{lmodern}
\ifPDFTeX\else
  % xetex/luatex font selection
  \setmainfont[]{Cambria}
  \setmathfont[]{Cambria Math}
\fi
% Use upquote if available, for straight quotes in verbatim environments
\IfFileExists{upquote.sty}{\usepackage{upquote}}{}
\IfFileExists{microtype.sty}{% use microtype if available
  \usepackage[]{microtype}
  \UseMicrotypeSet[protrusion]{basicmath} % disable protrusion for tt fonts
}{}
\makeatletter
\@ifundefined{KOMAClassName}{% if non-KOMA class
  \IfFileExists{parskip.sty}{%
    \usepackage{parskip}
  }{% else
    \setlength{\parindent}{0pt}
    \setlength{\parskip}{6pt plus 2pt minus 1pt}}
}{% if KOMA class
  \KOMAoptions{parskip=half}}
\makeatother
\usepackage{longtable,booktabs,array}
\newcounter{none} % for unnumbered tables
\usepackage{calc} % for calculating minipage widths
% Correct order of tables after \paragraph or \subparagraph
\usepackage{etoolbox}
\makeatletter
\patchcmd\longtable{\par}{\if@noskipsec\mbox{}\fi\par}{}{}
\makeatother
% Allow footnotes in longtable head/foot
\IfFileExists{footnotehyper.sty}{\usepackage{footnotehyper}}{\usepackage{footnote}}
\makesavenoteenv{longtable}
\setlength{\emergencystretch}{3em} % prevent overfull lines
\providecommand{\tightlist}{%
  \setlength{\itemsep}{0pt}\setlength{\parskip}{0pt}}
\usepackage{bookmark}
\IfFileExists{xurl.sty}{\usepackage{xurl}}{} % add URL line breaks if available
\urlstyle{same}
\hypersetup{
  pdftitle={The Roof and the Clock},
  pdfauthor={Selina Stenberg; Ilya Balashov},
  hidelinks,
  pdfcreator={LaTeX via pandoc}}

\title{The Roof and the Clock}
\usepackage{etoolbox}
\makeatletter
\providecommand{\subtitle}[1]{% add subtitle to \maketitle
  \apptocmd{\@title}{\par {\large #1 \par}}{}{}
}
\makeatother
\subtitle{Rowmotion on the E₇ minuscule poset, its linearity across the
minuscule family, and the Weyl group it does not share}
\author{Selina Stenberg \and Ilya Balashov}
\date{September 2026}

\begin{document}
\maketitle
\begin{abstract}
Let \(P\) be a simply-laced minuscule poset and \(R\) rowmotion on its
order ideals \(J(P)\), acting on the weights \(W\lambda\) of the
minuscule representation. By a theorem of Rush and Shi, \(R\) is
conjugate to a Coxeter element; a Coxeter element is an isometry of the
weights and \(R\) in general is not. We prove that \(R\) is itself a
Weyl group element exactly when \(P\) is a chain, that its half-turn
\(R^{h/2}\) is a Weyl group element exactly when \(P\) is a chain or
\(\lambda\) is the vector representation of \(D_n\), and we give in
closed form, on two families, the fraction of weight pairs whose inner
product \(R\) preserves. On the 56 weights of \(E_7\) a pair keeps its
inner product under one step of \(R\) if and only if a parity computed
from the two toggle sets is even. We then determine what a Weyl group
shares with its own translates by \(R\): \(W \cap R^kWR^{-k}\) equals
the centralizer \(C_W(R^k)\) on every board of rank at most 7 on which
no power of \(R\) is linear, with one exception, the \(D_5\) half-spin
(the \(B_4\) spinor board), whose mechanism we prove; on the spinor
boards \(B_4\), \(B_8\) and \(B_{16}\) the half-turn agrees with a Weyl
element of one fixed shape on the middle-rank orbits, and we prove that
this element commutes with the half-turn --- the property behind the
exception --- if and only if \(n = 4\). The longest element always
reverses \(R\), and three label statistics are homomesic. Every
computational statement is reproducible from the listed scripts.
\end{abstract}

\textbf{Keywords.} Rowmotion; minuscule posets; Coxeter elements; Weyl
groups; homomesy; the 56-dimensional representation of \(E_7\);
bitangents of plane quartics; the Fano plane.

\textbf{MSC 2020.} Primary 05E10; Secondary 06A07, 20F55, 05E18, 17B22.

\section{1. Introduction}\label{introduction}

\subsection{1.1 Rowmotion and the Coxeter
element}\label{rowmotion-and-the-coxeter-element}

For a finite poset \(P\), \emph{rowmotion} is the bijection of the set
\(J(P)\) of order ideals

\[R(I) = \text{the order ideal generated by the minimal elements of } P \setminus I .\]

Its study goes back to Brouwer--Schrijver, Fon-der-Flaass and
Cameron--Fon-der-Flaass {[}CF95{]}; the modern theory is surveyed in
{[}SW12, Hop24{]}. When \(P = P_\lambda\) is the \emph{minuscule poset}
of a minuscule weight \(\lambda\) of a simple Lie algebra --- the poset
of join-irreducibles of the distributive lattice formed by the weights
of \(V_\lambda\) under the root order {[}Pro84, Ste94, Gre13{]} --- Rush
and Shi proved that rowmotion is conjugate, under Stembridge's bijection
\(J(P_\lambda) \cong W/W_\lambda\) {[}Ste01{]}, to the action of any
Coxeter element \(c\) of the Weyl group \(W\) {[}RS13{]}. In particular
\(R\) has order the Coxeter number \(h\), its orbit structure is that of
\(c\), and \((J(P_\lambda), R)\) exhibits cyclic sieving {[}RSW04{]}.

A Coxeter element acts on the weights \(W\lambda\) as an isometry.
Rowmotion, carried to \(W\lambda\) by the same bijection, is a
permutation of the weights that is conjugate to \(c\) \emph{in the
symmetric group of the weights}, but it need not be an isometry, and in
general it is not even a Weyl group element. The first question of this
paper is: \textbf{for which minuscule posets is rowmotion, or a power of
rowmotion, itself an element of \(W\), and how far from an isometry is
it otherwise?}

\subsection{1.2 Results}\label{results}

Throughout, \(\lambda\) is a simply-laced minuscule weight (\(\omega_k\)
of \(A_n\); \(\omega_1, \omega_{n-1}, \omega_n\) of \(D_n\);
\(\omega_1, \omega_6\) of \(E_6\); \(\omega_7\) of \(E_7\)),
\(P = P_\lambda\), \(h\) the Coxeter number, \(W\) the Weyl group acting
on the weights \(W\lambda\), and \(R\) rowmotion on \(W\lambda\). Inner
products are normalised by \(\langle \alpha, \alpha \rangle = 2\) for
roots.

\textbf{Theorem A (linearity).} \emph{\(R \in W\) if and only if \(P\)
is a chain, i.e.~\(\lambda \in \{\omega_1, \omega_n\}\) of \(A_n\). If
\(h\) is even, \(R^{h/2} \in W\) if and only if \(P\) is a chain or
\(\lambda\) is the vector representation \(\omega_1\) of \(D_n\)
(\(n \ge 3\), with \(D_3\omega_1 = A_3\omega_2\) and, for \(D_4\), its
triality images); in the vector case \(R^{n-1}\) is the signed
permutation \(e_i \mapsto -e_{n-i}\) \((i<n)\),
\(e_n \mapsto (-1)^{n-1}e_n\). More generally, if \(R^j \in W\) then
\(j\) is a period of the extended rank sequence of \(P\) (§3).}

\textbf{Theorem B (the defect).} \emph{Let \(\kappa(R)\) be the fraction
of unordered pairs of weights whose inner product \(R\) preserves. Then
for \(n \ge 3\)}

\[\kappa(D_n\,\omega_1) = 1 - \frac{2(n-1)}{n(2n-1)}, \qquad \kappa(A_n\,\omega_2) = 1 - \frac{3n^2-9n+4}{\binom{\binom{n+1}{2}}{2}} .\]

\emph{Table 1 gives \(\kappa\), the random baseline, and the Hamming
distance from \(R\) to the nearest Weyl element for the 42 simply-laced
minuscule cases of rank \(\le 7\).}

\textbf{Theorem C (the parity rule on \(E_7\)).} \emph{Let \(\Lambda\)
be the 56 weights of the minuscule representation of \(E_7\), \(B\) the
inner product reduced mod 2, and for \(u \in \Lambda\) let \(T(u)\) be
the multiset of colours toggled by \(R\) at \(u\). For a non-antipodal
pair \(\{u,u'\}\) let \(t_1\) be the number of colours in \(T(u)\) at
which \(u'\) has a nonzero label, \(t_2\) the same with the roles
exchanged, and \(t_3\) the number of Dynkin-adjacent pairs between
\(T(u)\) and \(T(u')\). Then}

\[B(Ru, Ru') = B(u,u') \iff t_1 + t_2 + t_3 \equiv 0 \pmod 2 .\]

\emph{The rule holds with the \(k\)-fold toggle multisets for every
power \(R^k\). On \(E_7\), \(\kappa(R) = 501/770\), and no power of
\(R\) other than the identity is a Weyl element.}

The second question concerns what the symmetry group shares with the
clock. Write

\[I_k := W \cap R^k W R^{-k}, \qquad C_k := C_W(R^k),\]

so that \(I_k\) is the set of Weyl elements that are still Weyl elements
after conjugation by \(R^k\) (``the grammar two instants share'') and
\(C_k \subseteq I_k\).

\textbf{Theorem D (reversal).} \emph{On every simply-laced minuscule
board the longest element \(w_0\) reverses rowmotion:
\(w_0 R w_0 = R^{-1}\). The Weyl elements reversing \(R\) form the coset
\(w_0\,C_W(R)\); on \(E_7\) the reverser is unique and equals \(-1\). On
a self-dual board \(R^k\) preserves exactly
\(\tfrac12\,\mathrm{fix}(R^{2k})\) antipodal pairs.}

\textbf{Theorem E (the shared grammar).} \emph{(i) On \(E_7\),
\(I_k = \{1\}\) for \(k \ne 9\) and \(I_9 = \{\pm 1\}\); moreover
\(\langle W, R \rangle = A_{56}\). (ii) On the vector boards
\(D_n\,\omega_1\), \(I_k = C_W(R^{2k})\) for every \(k\); more
generally, if an involution \(\iota\) centralizes \(W\) and reverses
\(R\), then \(I_k \subseteq C_W(R^{2k})\). (iii) On the 23 boards of
rank \(\le 7\) on which no power of \(R\) is a Weyl element,
\(I_k = C_k\) for every \(k\), with one exception, the two half-spins of
\(D_5\) (equivalently the \(B_4\) spinor board, the vertices of the
4-cube), where \(|I_4| = 8 > |C_4| = 4\) and \(|I_3| = |I_5| = 2 > 1\).}

\textbf{Theorem F (the 4-cube).} \emph{On the \(B_4\) spinor board \(R\)
has two free orbits \(O_1 \ni \lambda\) and \(O_2\). The half-turn
\(R^4\) agrees with the Weyl element \(-\tau\),
\(\tau = (e_2 e_3)(-e_4)\), on all of \(O_2\) and with no Weyl element
on \(O_1\); consequently \(I_4 = C_W\big((-\tau)R^4\big)\), of order 8.
The obstruction on \(O_1\) is that \(O_1\) is the full-height orbit: for
any Weyl element \(w = (\pi,v)\) of the hyperoctahedral group the
rank-shift is
\(\mathrm{rank}(w x) - \mathrm{rank}(x) = \mathrm{wt}(u) - 2|x \cap u|\)
with \(u = \pi^{-1}v\), and an orbit containing \(\lambda\) and spanning
\(\mathbb{F}_2^n\) cannot carry a constant shift magnitude \(\ge 2\).}

The phenomenon does not recur on \(B_8\), the next board on which the
clock acts freely: there the half-turn again agrees with a Weyl element
on exactly two orbits, but the correction's centralizer is \(\{\pm 1\}\)
and \(I_k = C_W(R^k)\) at every lag (Proposition 6.7). The reason is a
commutation (Lemma 6.9): the Weyl element nearest to the half-turn
commutes with it on \(B_4\) and on no other enumerable spinor board
(Proposition 6.10); on \(B_{16}\), computed at the orbit level, the same
element again fails to commute and the correction's centralizer is again
\(\{\pm 1\}\) (Proposition 6.11). The last of these is a theorem for
every even \(n\):

\textbf{Theorem G (why four).} \emph{Let \(n \ge 4\) be even and let
\(w_n = -(e_2e_3)(e_4e_5)\cdots(e_{n-2}e_{n-1})\) be the signed
permutation negating \(e_1, \dots, e_{n-1}\), exchanging \(e_{2i}\) and
\(e_{2i+1}\), and fixing \(e_n\) (the notation is fixed in §6.4); for
\(n = 4, 8, 16\) this is the element that agrees with the half-turn
\(R^n\) on the middle-rank orbits of the \(B_n\) spinor board (Theorem
F, Propositions 6.7 and 6.11). Then \(w_n\) commutes with \(R^n\) if and
only if \(n = 4\). Consequently the correction \(w_n^{-1}R^n\) is an
involution only at \(n = 4\), and the mechanism of Theorem F cannot
enlarge the shared grammar on the agreement set of \(w_n\) for any even
\(n \ge 6\).} The proof (§6.4) rests on a closed form for rowmotion on
the shifted staircase in excess coordinates and for the orbit of
\(\lambda\), which also shows that the half-turn shifts every rank on
that orbit by exactly \(\pm n/2\) for every even \(n\); the two kinds of
orbit on which commutation is forced --- the orbit of \(\lambda\) and
the orbits carrying the half-turn --- exhaust the board exactly at
\(n = 4\), and for \(n \ge 6\) the weight with a single minus sign at
position 2 is an explicit witness.

We also show (§7) that the three ternary label counts on any
simply-laced minuscule board are homomesic under \(R\) with means
\(|P|/h\), \(|P|/h\), \(r - 2|P|/h\) --- Rush--Wang's antichain homomesy
{[}RW15{]} read on the Dynkin labels --- and that the toggle size
\(|I \,\triangle\, R(I)|\) is homomesic with mean \(2|P|/h\), a
corollary of Defant--Hopkins--Poznanović--Propp {[}DHPP23{]}.

The third part of the paper (§§8--9) is the explicit model in which
these objects were first met. Two finite groups built independently ---
a chain inside \(PSL(2,7)\) and a tower of double covers ---
\(SL(2,7)\), \(2\cdot I\), \(GL(2,3)\), \(2\cdot W(E_6)\),
\(2 \times PGL(2,7)\) --- are exhibited as subgroups of the rotation
subgroup \(W^+(E_8)\) over a single central \(C_2\); the quotient
\(O_8^+(2)\), its triality automorphism as an explicit permutation
\(\Phi\) of order 3 on 360 points, its fixed group \(G_2(2)\) as an
exact intersection, and the fourteen outer conjugacy classes of
\(O_8^+(2){:}3\) are all constructed in one coordinate system and
matched class by class to the ATLAS. The 56 weights of \(E_7\) are the
vectors \(u\) with \(B(u,v)=1\) for a fixed nonsingular \(v\); \(-1\) is
the transvection \(t_v\); and the involution

\[\mathrm{pr} := \iota \circ w_0(E_6) \circ t,\]

where \(t\) swaps the two ``poles'' of the \(\omega_7^\vee\)-grading, is
the order isomorphism between the two 27-sheets' ideal lattices induced
by the \(E_6\) diagram automorphism. Finally
\(\langle R, \mathrm{pr} \rangle = S_{56}\). Inside
\(Sp_6(2) \cong W(E_7)/\{\pm1\}\) acting on the 28 bitangents of a plane
quartic, \(W(E_6)\) is the stabilizer of a bitangent, \(PSL(2,7)\) is
transitive on all 28, and their intersection is \(S_3\); the 36 even
theta characteristics fall into \(PSL(2,7)\)-orbits of sizes
\(1, 7, 7, 21\), the two orbits of size 7 being the points and the lines
of one Fano plane, each realising the 6-dimensional irreducible
\(\chi_6\) as its deleted permutation module.

\subsection{1.3 What is classical, what is new, what is
computed}\label{what-is-classical-what-is-new-what-is-computed}

The Coxeter-number periodicity of \(R\), the structure of
\(O_8^+(2){:}3\), the fixed group of triality, the action of the \(E_6\)
diagram automorphism on minuscule posets, and the bitangent geometry of
\(Sp_6(2)\) are classical and are cited at first use. Theorems A--G and
Theorem C's rule are, to the best of a documented search (Appendix A),
not in the literature; we claim them as \emph{not found} rather than as
new, and would welcome a reference. Statements marked \textbf{(C)} were
established by exhaustive computation; the scripts, their check counts
and their input dependencies are listed in Appendix B, and each was
executed independently by both authors on separate machines with
identical tallies. Conjectures that we tested and found false are
recorded as Remarks where the failure is informative.

\subsection{1.4 Organisation}\label{organisation}

§2 fixes notation. §3 proves Theorems A and B. §4 treats the \(E_7\)
board: the parity rule, the shared grammar on the 56, and generation of
\(S_{56}\). §5 proves Theorem D and the vector-board part of Theorem E
and reports the family computation. §6 proves Theorem F, describes the
chain \(P_j(k)\) of iterated shared grammars, reports the \(B_8\) and
\(B_{16}\) computations, and proves Theorem G. §7 is homomesy. §8
describes the computational model of \(W^+(E_8)\), \(O_8^+(2){:}3\) and
\(G_2(2)\) and identifies the 56, the transvection and the mirror. §9 is
the bitangent bridge and the Fano plane. §10 lists open problems.
Appendix A is the literature search; Appendix B the verification table.

\section{2. Preliminaries}\label{preliminaries}

\textbf{Weights and posets.} For a minuscule weight \(\lambda\) the
weights of \(V_\lambda\) form the single orbit \(W\lambda\); each weight
\(w\) has Dynkin labels
\(\langle w, \alpha_i^\vee \rangle \in \{-1,0,1\}\). The root order
\(w \le w'\) iff \(w' - w \in \mathbb{N}\{\alpha_i\}\) makes
\(W\lambda\) a distributive lattice with top \(\lambda\) and bottom
\(w_0\lambda = -\lambda^*\); its join-irreducibles form the minuscule
poset \(P = P_\lambda\), coloured by simple roots, and the ideals
\(J(P)\) correspond to the weights, \(\lambda \leftrightarrow P\),
\(w_0\lambda \leftrightarrow \varnothing\), with \$w(I) - w(I') = \$ the
sum of the colours of \(I' \setminus I\) for \(I \subseteq I'\)
{[}Pro84, Ste94{]}. The multiplicity \(c_i\) of colour \(i\) in \(P\) is
the coefficient of \(\alpha_i\) in
\(\lambda - w_0\lambda = \lambda + \lambda^*\). \(P\) is graded; we
write \(s_j\) for the number of elements of rank \(j\),
\(0 \le j \le h-2\).

\textbf{Rowmotion on weights.} Under \(J(P) \cong W\lambda\), \(R\)
sends the weight of \(I\) to the weight of
\(\downarrow\!\min(P\setminus I)\). The multiset \(T(w)\) of colours of
the elements of \(I \,\triangle\, R(I)\) is the \emph{toggle set} at
\(w\); the weight changes by the signed sum of the corresponding simple
roots. \textbf{Theorem (Rush--Shi {[}RS13{]}).} \(R\) is conjugate in
the toggle group to the product of the toggles in any linear order, and
under Stembridge's bijection that product corresponds to a Coxeter
element; hence \(R\) has order \(h\) and the cycle type of a Coxeter
element on \(W\lambda\).

\textbf{Boards.} We call the pair \((W \curvearrowright W\lambda, R)\) a
\emph{board}; it is \emph{rich} when no power of \(R\) is a Weyl element
and \(C_W(R) = 1\). A board is \emph{self-dual} when \(-1\) permutes
\(W\lambda\) (i.e.~\(\lambda^* = \lambda\)); then the \emph{antipode}
\(\iota: w \mapsto -w\) is a permutation of the weights, equal to
\(w_0\) when \(-1 \in W\) (\(D_n\), \(n\) even; \(E_7\)) and not a Weyl
element otherwise (\(A_{2k-1}\omega_k\), \(D_n\omega_1\) with \(n\)
odd).

\textbf{The hyperoctahedral case.} The spinor representation of \(B_n\)
is minuscule with weights \((\pm\tfrac12)^n\); we read a weight as
\(x \in \mathbb{F}_2^n\) (bit \(i = 1\) iff coordinate \(i\) is
negative), so \(\mathrm{rank}(x) = \mathrm{wt}(x)\), \(\lambda = 0\),
\(w_0\lambda = \mathbf{1}\). Its poset is the shifted staircase
\(\delta_n\) and \(W(B_n) = \mathbb{F}_2^n \rtimes S_n\) acts affinely,
\(w(x) = \pi(x) \oplus v\). The \(D_{n+1}\) half-spin has the same poset
and the same rowmotion, with \(W(B_n) < W(D_{n+1})\).

\textbf{Notation for the shared grammar.} For a board and
\(1 \le k \le h-1\), \(I_k = W \cap R^kWR^{-k}\) and \(C_k = C_W(R^k)\);
the \emph{transport} is the map \(g \mapsto R^{-k} g R^{k}\) on \(I_k\)
(which by definition lands in \(W\)). Conjugating by \(R^k\) shows
\(|I_k| = |I_{-k}|\).

\section{3. Linearity of rowmotion}\label{linearity-of-rowmotion}

Let \(\sigma = (s_0, s_1, \dots, s_{h-2}, c_k)\) for
\(\lambda = \omega_k\); we call \(\sigma\), of length \(h\), the
\emph{extended rank sequence}, indices taken mod \(h\).

\textbf{Lemma 3.1.} \emph{(a) \(R(\lambda) = w_0\lambda\); the bottom
weight \(w_0\lambda\) has a single nonzero label, \(-1\) at the dual
node \(k^*\), hence a single cover \(p_0 = w_0\lambda + \alpha_{k^*}\),
the unique minimal element of \(P\), and \(R(w_0\lambda) = p_0\).
Consequently}

\[\langle \lambda, w_0\lambda \rangle = \langle \lambda, \lambda \rangle - c_k, \qquad \langle w_0\lambda, p_0 \rangle = \langle \lambda, \lambda \rangle - 1 .\]

\emph{(b) The \(R\)-orbit of \(w_0\lambda\) is the sequence of rank
truncations \(P_{<j}\), \(j = 0, \dots, h-1\); writing \(w(j)\) for
their weights,}

\[\langle w(j), w(j+1) \rangle = \langle \lambda, \lambda \rangle - \sigma_j \quad (j \bmod h).\]

\emph{Proof.} (a) \(R(P) = \downarrow\!\min(\varnothing) = \varnothing\)
and \(R(\varnothing) = \downarrow\!\min P = \{p_0\}\). Since
\(\lambda - w_0\lambda = \sum_i c_i\alpha_i\) and
\(\langle \omega_k, \alpha_i \rangle = \delta_{ik}\),
\(\langle \lambda, \lambda - w_0\lambda \rangle = c_k\); and
\(\langle w_0\lambda, \alpha_{k^*} \rangle = -1\) because that is the
label of \(w_0\lambda\) at \(k^*\). (b) The minimal elements of
\(P \setminus P_{<j}\) are the rank-\(j\) elements, which generate
\(P_{<j+1}\); so \(w(j+1) - w(j)\) is the sum of the colours of the
rank-\(j\) elements, each added at a label \(-1\) of \(w(j)\), and the
inner product drops by \(s_j\); all weights have the same norm. The
wrap-around step \(\lambda \to w_0\lambda\) drops by \(c_k\) by (a).
\(\square\)

\textbf{Theorem 3.2 (= Theorem A, first part).} \emph{\(R \in W \iff P\)
is a chain \(\iff \lambda \in \{\omega_1, \omega_n\}\) of \(A_n\); then
\(R\) is the Coxeter element \(s_1 \cdots s_n\), the cyclic shift of
\(e_1, \dots, e_{n+1}\).}

\emph{Proof.} If \(R \in W\) it is an isometry, so applying it to the
pair \((\lambda, w_0\lambda) \mapsto (w_0\lambda, p_0)\) of Lemma 3.1(a)
gives
\(\langle\lambda,\lambda\rangle - c_k = \langle\lambda,\lambda\rangle - 1\),
i.e.~\(c_k = 1\). From the classification of minuscule posets,
\(c_k = \min(k, n+1-k)\) for \(A_n\omega_k\), \(2\) for \(D_n\omega_1\),
\(\lfloor n/2 \rfloor\) for \(D_n\omega_n\) (\(n \ge 4\)), \(2\) for
\(E_6\omega_1\) and \(3\) for \(E_7\omega_7\); so \(c_k = 1\) only for
\(A_n\) with \(k \in \{1,n\}\), whose poset is a chain. Conversely on a
chain the ideals are the initial segments and \(R\) is the cyclic shift,
a permutation matrix in \(W(A_n) = S_{n+1}\); that this shift is the
Coxeter element \(s_1 \cdots s_n\) is classical (cf.~{[}SW12, Thm
6.1{]}, attributing the chain and product-of-chains cases to Stanley).
\(\square\)

\textbf{Theorem 3.3 (= Theorem A, second part).} \emph{Let \(h\) be even
and \(m = h/2\). Then \(R^m \in W\) if and only if \(P\) is a chain or
\(\lambda\) is the vector representation \(\omega_1\) of \(D_n\),
\(n \ge 3\). In the vector case, on the weights \(\pm e_i\),}

\[R^{n-1}(e_i) = -e_{n-i} \ (1 \le i \le n-1), \qquad R^{n-1}(e_n) = (-1)^{n-1} e_n,\]

\emph{a signed permutation with an even number of sign changes, hence in
\(W(D_n)\).}

\emph{Proof.} (Obstruction.) If \(R^m \in W\) it is an isometry with
\(R^m(w(j)) = w(j+m)\), so Lemma 3.1(b) gives
\(\sigma_{j+m} = \sigma_j\) for all \(j\): \(\sigma\) is \(m\)-periodic,
and since \(\sigma_0 = s_0 = 1\), necessarily \(\sigma_m = 1\). For
\(A_n\omega_k\) with \(1 < k < n\) (\(n\) odd, \(m = (n+1)/2\)), \(P\)
is the \(k \times (n+1-k)\) rectangle and
\(s_m = \min(k, n+1-k) \ge 2\). For \(D_n\omega_n\), \(n \ge 5\)
(\(m = n-1\)), \(P\) is the shifted staircase
\(\{(i,j) : 1 \le i \le j \le n-1\}\) with rank \(i+j-2\), and
\(s_{n-1} = \lfloor (n+1)/2 \rfloor - 1 \ge 2\). For \(E_6\omega_1\)
(\(m = 6\)) the rank sizes are \(1,1,1,2,2,2,2,2,1,1,1\) and
\(s_6 = 2\); for \(E_7\omega_7\) (\(m = 9\)) they are
\(1,1,1,1,2,2,2,2,3,2,2,2,2,1,1,1,1\) and \(s_9 = 2\). For
\(D_n\omega_1\) (\(m = n-1\)), \(P\) is a chain of \(n-2\) elements
above the antichain \(\{e_n, -e_n\}\) above a chain of \(n-2\), with
\(c_1 = 2\), and \(\sigma = (1^{n-2}, 2, 1^{n-2}, 2)\) is
\((n-1)\)-periodic; the triality images at \(n = 4\) and the chains pass
likewise.

(Sufficiency.) On \(D_n\omega_1\) the ideals are \(\varnothing\), the
prefixes of the lower chain, the three ideals adding \(e_n\), \(-e_n\)
or both, and the extensions along the upper chain; rowmotion cycles
\(\varnothing \to \cdots \to P\) in \(2n-2\) steps and swaps the two
one-sided ideals at every step. On weights,

\[-e_1 \to -e_2 \to \cdots \to -e_{n-1} \to e_{n-1} \to \cdots \to e_1 \to -e_1, \qquad e_n \leftrightarrow -e_n .\]

The cycle has length \(2n-2\) with \(-e_i\) at position \(i-1\) and
\(e_j\) at position \(2n-2-j\); shifting by \(n-1\) sends \(-e_i\) to
position \(n+i-2 = 2n-2-(n-i)\), i.e.~to \(e_{n-i}\). Hence
\(R^{n-1}(\pm e_i) = \mp e_{n-i}\) for \(i \le n-1\) and
\(R^{n-1}(e_n) = (-1)^{n-1}e_n\): the coordinate permutation
\(i \leftrightarrow n-i\) with \(n-1\) sign changes on the first \(n-1\)
coordinates and one more iff \(n\) is even --- \(n\) or \(n-1\) sign
changes, even either way. \(\square\)

\textbf{Corollary 3.4.} \emph{\(R^j \in W\) implies that \(\sigma\) is
\(j\)-periodic; so the linear powers of rowmotion are among the periods
of \(\sigma\). For \(E_7\), \(\sigma\) has no period below \(18\): no
power of \(R\) but the identity is a Weyl element.}

\textbf{Theorem 3.5 (= Theorem B).} \emph{Let \(\kappa(R)\) be the
fraction of unordered pairs \(\{w,w'\}\) of weights with
\(\langle Rw, Rw' \rangle = \langle w, w' \rangle\). For \(n \ge 3\),}

\[\kappa(D_n\,\omega_1) = 1 - \frac{2(n-1)}{n(2n-1)}, \qquad \kappa(A_n\,\omega_2) = 1 - \frac{3n^2 - 9n + 4}{\binom{\binom{n+1}{2}}{2}} .\]

\emph{Proof.} (\(D_n\omega_1\).) The inner products are \(0\) and
\(-1\), the latter on the \(n\) antipodal pairs. By the cycle in the
proof of Theorem 3.3, \(\{e_n, -e_n\}\) is the only antipodal pair \(R\)
keeps; the other \(n-1\) become non-antipodal, and since \(R\) permutes
the weights bijectively, \(n-1\) non-antipodal pairs become antipodal:
\(2(n-1)\) of the \(n(2n-1)\) pairs change.

(\(A_n\omega_2\).) Weights are the 2-subsets \(\{a < b\}\) of \([n+1]\)
and a pair's inner product is determined by \(|S \cap T| \in \{0,1\}\).
The ideals \((x_1 \ge x_2)\) of the \(2 \times (n-1)\) rectangle
correspond to \(\{n - x_1, n+1-x_2\}\), and rowmotion on the rectangle
gives

\[R\{a,b\} = \{a-1,\, b-1\} \ (a \ge 2,\ b \ge a+2), \quad R\{a,a+1\} = \{a-1,\, n+1\} \ (a \ge 2),\]

\[R\{1,b\} = \{b-2,\, b-1\} \ (b \ge 3), \quad R\{1,2\} = \{n,\, n+1\} .\]

Let \(\rho\) be the rotation \(i \mapsto i-1 \pmod{n+1}\), a Coxeter
element and an isometry. Then \(R = \rho \circ \varphi\), where
\(\varphi\) fixes every 2-subset outside
\(T = \{\{a,a+1\} : 2 \le a \le n\} \cup \{\{1,b\} : 2 \le b \le n+1\}\)
and permutes \(T\) by \(\{a,a+1\} \mapsto \{1,a\}\),
\(\{1,b\} \mapsto \{b-1,b\}\) (\(b \ge 3\)),
\(\{1,2\} \mapsto \{1,n+1\}\); so \(\kappa(R) = \kappa(\varphi)\). A
pair \((S \in T, G \notin T)\) changes type iff \(a+1 \in G\) (for
\(S = \{a,a+1\}\)), iff \(b-1 \in G\) (for \(S = \{1,b\}\),
\(b \ge 3\)), or iff exactly one of \(2, n+1\) lies in \(G\) (for
\(S = \{1,2\}\)): \(2(n-2)^2 + 2(n-3)\) pairs. Inside \(T\), of the
\((n-2) + 2(n-1) + \binom{n}{2}\) intersecting pairs exactly \(4n-5\)
remain intersecting, so \((n-1)(n-2)\) pairs change. The total is
\(3n^2 - 9n + 4\). \(\square\)

Every step of these proofs was re-checked on the data \textbf{(C)}:
Lemma 3.1 on all 42 cases of Table 1; \(c_k = 1 \iff\) chain
\(\iff R \in W\); the periodicity criterion with exactly the predicted
case list; the vector formula for \(n = 3, \dots, 9\); the explicit
rowmotion on 2-subsets and both formulas of Theorem 3.5 for
\(n = 3, \dots, 9\).

\textbf{Remark 3.6 (observed).} The analogous count for \(A_n\omega_3\)
is, for \(n \le 10\), the degree-4 polynomial
\(12 D_3(n) = 15n^4 - 110n^3 + 165n^2 + 314n - 432\), suggesting degree
\(2k-2\) for \(\omega_k\). We have no proof.

\subsection{3.1 The defect across the
family}\label{the-defect-across-the-family}

One program built from the Cartan matrix alone (weights in Dynkin
labels, the lattice by \(w \mapsto w - \alpha_i\) at label \(1\), its
join-irreducibles, rowmotion as the join of the minimal elements of the
complement) computes, for the 42 simply-laced minuscule cases with
\(A_n\) (\(2 \le n \le 7\), all \(k\); the trivial \(A_1\) omitted),
\(D_n\) (\(n = 4..7\), \(\omega_1\) and both half-spins), \(E_6\) and
\(E_7\): \(\kappa\); the random baseline \(\kappa_0 = \sum_t p_t^2\)
over the distribution \(p_t\) of inner-product types; the best pointwise
agreement \(a(R)\) of \(R\) with a Coxeter element; the \emph{altitude}
\(\nu(R)\), the Hamming distance from \(R\) to the nearest element of
\(W\) acting on the weights (\(W\) enumerated as permutations of the
weights where \(|W| \le 322{,}560\); for \(E_7\) by a
base-and-strong-generating-set search over all \(2{,}903{,}040\)
elements); and the linear centralizer \(C_W(R)\).

\textbf{Table 1.} Rowmotion against the Weyl group on the simply-laced
minuscule boards of rank \(\le 7\) (dual representations give equal
rows).

{\def\LTcaptype{none} % do not increment counter
\begin{longtable}[]{@{}
  >{\raggedright\arraybackslash}p{(\linewidth - 16\tabcolsep) * \real{0.1111}}
  >{\raggedright\arraybackslash}p{(\linewidth - 16\tabcolsep) * \real{0.1111}}
  >{\raggedright\arraybackslash}p{(\linewidth - 16\tabcolsep) * \real{0.1111}}
  >{\raggedright\arraybackslash}p{(\linewidth - 16\tabcolsep) * \real{0.1111}}
  >{\raggedright\arraybackslash}p{(\linewidth - 16\tabcolsep) * \real{0.1111}}
  >{\raggedright\arraybackslash}p{(\linewidth - 16\tabcolsep) * \real{0.1111}}
  >{\raggedright\arraybackslash}p{(\linewidth - 16\tabcolsep) * \real{0.1111}}
  >{\raggedright\arraybackslash}p{(\linewidth - 16\tabcolsep) * \real{0.1111}}
  >{\raggedright\arraybackslash}p{(\linewidth - 16\tabcolsep) * \real{0.1111}}@{}}
\toprule\noalign{}
\begin{minipage}[b]{\linewidth}\raggedright
case
\end{minipage} & \begin{minipage}[b]{\linewidth}\raggedright
\(|W\lambda|\)
\end{minipage} & \begin{minipage}[b]{\linewidth}\raggedright
\(h\)
\end{minipage} & \begin{minipage}[b]{\linewidth}\raggedright
cycle type of \(R\)
\end{minipage} & \begin{minipage}[b]{\linewidth}\raggedright
\(\kappa(R)\)
\end{minipage} & \begin{minipage}[b]{\linewidth}\raggedright
\(\kappa_0\)
\end{minipage} & \begin{minipage}[b]{\linewidth}\raggedright
\(a(R)\)
\end{minipage} & \begin{minipage}[b]{\linewidth}\raggedright
\(\nu(R)\)
\end{minipage} & \begin{minipage}[b]{\linewidth}\raggedright
\(|C_W(R)|\)
\end{minipage} \\
\midrule\noalign{}
\endhead
\bottomrule\noalign{}
\endlastfoot
\(A_n\,\omega_1\) & \(n+1\) & \(n+1\) & \((n+1)\) & \(1\) & \(1\) &
\(n+1\) & \(0\) & \(n+1\) \\
\(A_3\,\omega_2 = D_3\,\omega_1\) & 6 & 4 & \(4\cdot 2\) & \(11/15\) &
0.680 & 3 & 3 & 2 \\
\(A_4\,\omega_2\) & 10 & 5 & \(5^2\) & \(29/45\) & 0.556 & 5 & 5 & 1 \\
\(A_5\,\omega_2\) & 15 & 6 & \(6^2\cdot 3\) & \(71/105\) & 0.510 & 6 & 9
& 1 \\
\(A_5\,\omega_3\) & 20 & 6 & \(6^3\cdot 2\) & \(58/95\) & 0.452 & 7 & 13
& 1 \\
\(A_6\,\omega_2\) & 21 & 7 & \(7^3\) & \(76/105\) & 0.500 & 10 & 11 &
1 \\
\(A_6\,\omega_3\) & 35 & 7 & \(7^5\) & \(339/595\) & 0.419 & 10 & 25 &
1 \\
\(A_7\,\omega_2\) & 28 & 8 & \(8^3\cdot 4\) & \(145/189\) & 0.506 & 15 &
13 & 1 \\
\(A_7\,\omega_3\) & 56 & 8 & \(8^7\) & \(431/770\) & 0.405 & 12 & 44 &
1 \\
\(A_7\,\omega_4\) & 70 & 8 & \(8^8\cdot 4\cdot 2\) & \(417/805\) & 0.380
& 12 & 56 & 1 \\
\(D_4\,\omega_1\) & 8 & 6 & \(6\cdot 2\) & \(11/14\) & 0.755 & 4 & 3 &
2 \\
\(D_5\,\omega_1\) & 10 & 8 & \(8\cdot 2\) & \(37/45\) & 0.802 & 5 & 5 &
2 \\
\(D_5\,\omega_5\) & 16 & 8 & \(8^2\) & \(19/30\) & 0.556 & 5 & 10 & 1 \\
\(D_6\,\omega_1\) & 12 & 10 & \(10\cdot 2\) & \(28/33\) & 0.835 & 6 & 5
& 2 \\
\(D_6\,\omega_6\) & 32 & 10 & \(10^3\cdot 2\) & \(79/124\) & 0.469 & 9 &
19 & 1 \\
\(D_7\,\omega_1\) & 14 & 12 & \(12\cdot 2\) & \(79/91\) & 0.858 & 7 & 7
& 2 \\
\(D_7\,\omega_7\) & 64 & 12 & \(12^5\cdot 4\) & \(295/504\) & 0.432 & 13
& 48 & 1 \\
\(E_6\,\omega_1\) & 27 & 12 & \(12^2\cdot 3\) & \(223/351\) & 0.527 & 9
& 15 & 1 \\
\(E_7\,\omega_7\) & 56 & 18 & \(18^3\cdot 2\) & \(501/770\) & 0.482 & 14
& 38 & 1 \\
\end{longtable}
}

Rush--Shi's conjugacy is visible in every row (the cycle type is that of
a Coxeter element). \(\kappa\) exceeds \(\kappa_0\) in every non-chain
case. The linear centralizer is trivial except on chains and on the
vector representations, where it is \(\{1, R^{h/2}\}\) (Theorem 3.3).
\textbf{Remark 3.7.} \(\kappa\) is not monotone in the rank along every
family: it increases along \(D_n\omega_1\) and \(A_n\omega_2\) (Theorem
3.5) and decreases along the half-spins. \(E_6\) and \(E_7\) sit where
the \(D_5\)--\(D_7\) spinors sit; the \(E_7\) board is, by every measure
in the table, as far from linear as rowmotion on a minuscule poset gets
at this rank.

\section{\texorpdfstring{4. The \(E_7\)
board}{4. The E\_7 board}}\label{the-e_7-board}

Let \(\Lambda = W\lambda\) be the 56 weights of the minuscule
representation of \(E_7\), \(W = W(E_7)\) of order
\(2{,}903{,}040 = 2 \times |Sp_6(2)|\), \(P\) the 27-element minuscule
poset, and \(R\) rowmotion, of order 18 with cycle type \(18^3\cdot 2\).
Since \(-1 \in W(E_7)\), \(\iota := -1 = w_0\) is central in \(W\) and
pairs the weights into 28 antipodal pairs. With
\(\langle \alpha,\alpha\rangle = 2\) the weights have norm \(3/2\) and
two distinct weights have inner product \(1/2\), \(-1/2\) or \(-3/2\),
the last exactly for antipodal pairs; each weight has 27 partners of
each non-antipodal type. Write \(B(u,u') \in \mathbb{F}_2\) for the
\emph{type} of a non-antipodal pair, \(0\) if
\(\langle u,u' \rangle = 1/2\) and \(1\) if \(-1/2\). \(B\) is the
restriction to \(\Lambda\) of the symmetric bilinear form of the
\(\mathbb{F}_2\)-model of §8.3, in which the 56 weights are the
nonsingular vectors \(u\) with \(B(u,v)=1\); in particular \(B\) is
bilinear on the \(\mathbb{F}_2\)-span of \(\Lambda\), and
\(B(\beta, w) = \langle \alpha, w \rangle \bmod 2\) for a simple root
\(\alpha\) with image \(\beta\).

\subsection{4.1 The parity rule}\label{the-parity-rule}

For \(u \in \Lambda\) write \(\delta(u) := u + R(u)\) (in
\(\mathbb{F}_2\)-coordinates); \(\delta(u)\) is the sum of the simple
roots \(\beta_c\), \(c \in T(u)\), toggled at \(u\). The toggle sizes on
the 56 are
\(\{1 : 8,\ 2 : 28,\ 3 : 11,\ 4 : 2,\ 5 : 3,\ 6 : 2,\ 9 : 1,\ 27 : 1\}\).

\textbf{Proposition 4.1 (= Theorem C).} \emph{For a non-antipodal pair
\(\{u,u'\}\), \(B(Ru, Ru') = B(u,u')\) iff \(t_1 + t_2 + t_3\) is even,
where \(t_1 = \#\{c \in T(u) : \  ext{label}_c(u') \ne 0\}\), \(t_2\)
likewise with the roles exchanged, and \(t_3\) is the number of
Dynkin-adjacent pairs \((c, c') \in T(u) \times T(u')\).}

\emph{Proof.} \(B\) is bilinear, so

\[B(Ru, Ru') - B(u,u') = B(\delta u, u') + B(u, \delta u') + B(\delta u, \delta u') .\]

With \(\delta u = \sum_{c \in T(u)} \beta_c\) we have
\(B(\beta_c, w') = \langle \alpha_c, w' \rangle \bmod 2 = |\ ext{label}_c(w')|\)
for minuscule labels in \(\{-1,0,1\}\), and \(B(\beta_c, \beta_{c'})\)
is the Cartan entry mod 2, i.e.~Dynkin adjacency. The three terms are
\(t_1, t_2, t_3\). \(\square\)

The content of the proposition is therefore the \emph{count}: on all
\(1{,}485\) non-antipodal pairs the rule has zero exceptions
\textbf{(C)}, and with the 28 antipodal pairs the kept count is
\(1{,}002\) of \(1{,}540\), i.e.~\(\kappa(R) = 501/770\) as in Table 1.
The eight cells \((t_1,t_2,t_3)\) and their populations are tabulated in
the verification log; the dominant kept cell is \((1,1,0)\) with 553
pairs. The same rule with the \(k\)-fold toggle multiset is exact for
every power \(R^k\), \(k = 1, \dots, 17\), and reproduces the measured
profile \(\kappa(R^9) = 0.616\), \(\kappa(R^2) = 0.547\),
\(\kappa(R^3) = 0.513\) against a random mean of \(0.482\). Each
27-sheet (colour 7 deleted) is \(J(P_6)\) with \(|P_6| = 16\), its
rowmotion \(R_6\) has order \(12\) with orbits \(12, 12, 3\), its pair
types are the Schläfli adjacency, and the same rule with the \(E_6\)
diagram is exact there; \(\kappa(R_6) = 223/351 = 0.635 < \kappa(R)\)
--- smaller toggle sets do not mean fewer flips.

\subsection{4.2 What the Weyl group shares with the clock on the
56}\label{what-the-weyl-group-shares-with-the-clock-on-the-56}

\textbf{Proposition 4.2 (C).} \emph{\(C_{S_{56}}(R) \cap W = \{1\}\): no
non-trivial Weyl element commutes with \(R\). More strongly, \(R\)
agrees pointwise with the nearest Weyl element on 18 of the 56 weights,
so \(\nu(R) = 38\), and with the nearest Coxeter element on 14 (Table
1); the powers have altitudes \(\nu(R^9) = 36\),
\(\nu(R^2) = \nu(R^6) = 44\), \(\nu(R^3) = 46\).}

\textbf{Theorem 4.3 (= Theorem E(i)) (C).}

\[W \cap R^kWR^{-k} = \{1\} \quad (k = 1,\dots,8,\ 10,\dots,17), \qquad W \cap R^9WR^{-9} = \{1, \iota\} ;\]

\[\langle W, R \rangle = A_{56}, \qquad \langle W, \mathrm{pr} \rangle = 2^{28} \rtimes Sp_6(2), \qquad \langle W, R, \mathrm{pr} \rangle = S_{56},\]

\emph{where \(\mathrm{pr}\) is the involution of §4.3.}

\emph{Proof (C).} \(I_k\) was computed by conjugating every element of
\(W\) by \(R^k\) and testing membership in \(W\) (a permutation of
\(\Lambda\) is in \(W\) iff it extends to an isometry of the weight
configuration; for \(E_7\) this criterion is exact because the isometry
group of \(\Lambda\) is \(W\) itself); for \(k \ne 9\) the cheap
necessary test --- the conjugate must commute with the central \(\iota\)
--- already leaves only the identity. That \(\iota \in I_9\) is forced:
\(R^9\) commutes with \(\iota\) (Theorem D) and \(\iota\) is central in
\(W\). \(R\), \(\iota\) and all reflections are even permutations of the
56 and \(\mathrm{pr}\) is odd; the group orders were established by
Schreier--Sims. For \(\langle W, \mathrm{pr} \rangle\): \(\mathrm{pr}\)
fixes exactly the two poles, which form an antipodal pair, and
\(\mathrm{pr} \circ t\) (\(t\) the pole swap) is a Weyl element (§4.3);
\(W\) is transitive on the 28 antipodal pairs with image \(Sp_6(2)\) and
kernel \(\langle\iota\rangle\), so the normal closure of \(t\) in
\(\langle W, \mathrm{pr} \rangle\) is the group \(2^{28}\) of all
pair-flips, meeting \(W\) in \(\langle \iota \rangle\). \(\square\)

Corollaries: no reflection of \(W\) (0 of 63) is a symmetry after any
tick \(k = 1, \dots, 17\), and none of the 21 involutions of a bridge
subgroup \(PSL(2,7) \le W\) (§9) is; the eighteen conjugates
\(R^kWR^{-k}\) are eighteen copies of \(W(E_7)\) in \(S_{56}\) meeting
pairwise in the identity (in \(\langle\iota\rangle\) at distance 9).

\subsection{4.3 The mirror}\label{the-mirror}

The coweight \(\omega_7^\vee\) takes the values \(\pm\tfrac32\) on two
weights (the \emph{poles}) and \(\pm\tfrac12\) on two sets of 27 (the
\emph{sheets}); each sheet is the weight lattice of a 27-dimensional
representation of \(E_6\), with its own rowmotion \(R_6\) and its own
longest element \(w_0(E_6) \in W(E_6) < W(E_7)\), acting on the sheet's
\(E_6\)-labels as \(-\sigma_6\), where \(\sigma_6\) is the \(E_6\)
diagram automorphism. Define

\[\mathrm{pr} := \iota \circ w_0(E_6) \circ t, \qquad t = \text{the transposition of the two poles}.\]

\textbf{Proposition 4.4 (C).} \emph{(i) \(\mathrm{pr}\) is an involution
swapping the two sheets and fixing the poles; on the sheets it is the
order isomorphism \(J(P_6^{(0)}) \to J(P_6^{(1)})\) whose induced map on
join-irreducibles is the colour map \(0 \leftrightarrow 5\),
\(2 \leftrightarrow 4\), \(1, 3\) fixed --- the \(E_6\) diagram
automorphism (324 of 324 comparable pairs preserved each way, principal
ideals to principal ideals). (ii) \(\mathrm{pr}\) commutes with \(R_6\)
and with \(\iota\), and \(\mathrm{pr} \circ t = \iota \circ w_0(E_6)\)
is a fixed-point-free involution of \(W(E_7)\) at Hamming distance 2
from \(\mathrm{pr}\) (so \(\nu(\mathrm{pr}) = 2\)). (iii)
\(\mathrm{pr}\) preserves the inner-product type of exactly \(1{,}432\)
of the \(1{,}540\) pairs --- the count of an isometry composed with one
transposition. (iv) \(|\langle R, \iota \rangle| = 36\),
\(|\langle \mathrm{pr}, \iota \rangle| = 4\), and
\(\langle R, \mathrm{pr} \rangle = S_{56}\).}

That the diagram automorphism swaps the two 27's and induces an
isomorphism of their minuscule posets is standard {[}Gre13; GX26{]};
(ii) follows from Theorem D applied on each sheet (\(w_0(E_6)\) reverses
\(R_6\), \(\iota\) reverses both \(R\) and \(R_6\), and a product of two
reversers commutes). (iv) was established by Schreier--Sims; the Cayley
diameter of \(\{R, \mathrm{pr}\}\) in \(S_{56}\) is at least 157 by
counting.

\textbf{Proposition 4.5 (the reversers on the 56) (C).} \emph{Exactly
one element of \(W(E_7)\) reverses \(R\), namely \(\iota\); exactly one
element of \(W(E_6)\) reverses \(R_6\), namely \(w_0(E_6)\). Hence
\(\mathrm{pr}\) is the product of the two reversals, corrected at the
poles.}

\subsection{4.4 The bridge subgroup and the outer
automorphism}\label{the-bridge-subgroup-and-the-outer-automorphism}

Let \(L \le W(E_7)\) be a copy of \(PSL(2,7)\) transitive on the 28
antipodal pairs (§9 constructs it; its normaliser in \(Sp_6(2)\) is
computed in Proposition 9.1(iv)).

\textbf{Proposition 4.6 (C).}
\emph{\(C_{W(E_7)}(L) = \langle \iota \rangle\) and \(N_{W(E_7)}(L)\)
has order 672, so \(N/C \cong PGL(2,7)\): \(W(E_7)\) realises every
automorphism of \(L\) and no further symmetry of it.
\(N = \langle\iota\rangle \times PGL(2,7)\), with one \(PGL(2,7)\)
complement inside \(Sp_6(2) = W^+\) (fixed-point-free on \(\Lambda\))
and one outside. The outer coset contains 56 involutions (28 of cycle
type \(2^{24}1^8\), 28 fixed-point-free), 112 elements of order 6 and
168 of order 8. Any \(c \in N \setminus L\langle\iota\rangle\) is a Weyl
element and therefore commutes with no power of \(R\) and reverses none
(Proposition 4.2, Theorem 4.3): the unique reverser of \(R\) in \(N\) is
\(\iota\) and the unique element of \(N\) commuting with \(R\) is
\(1\).}

\emph{\(P\) acts on the 28 antipodal pairs with permutation character
\(\chi_1 + 2\chi_6 + \chi_7 + \chi_8\) (fixed points \(28, 4, 1, 0, 0\)
on the classes \(1A, 2A, 3A, 4A, 7AB\)), and
\(\mathbb{C}\Lambda = 2(\chi_1 + 2\chi_6 + \chi_7 + \chi_8)\) as a
\(P\)-module; the three-dimensional irreducibles \(\chi_3, \bar\chi_3\)
do not occur.}

\section{5. The shared grammar across the
family}\label{the-shared-grammar-across-the-family}

\subsection{5.1 Reversal}\label{reversal}

\textbf{Theorem 5.1 (= Theorem D).} \emph{On every simply-laced
minuscule board the longest element \(w_0\), acting on weights as
\(-\sigma\) (\(\sigma\) the diagram automorphism carrying \(\lambda\) to
\(\lambda^*\); trivial when \(-1 \in W\)), reverses rowmotion:}

\[w_0 R w_0 = R^{-1} .\]

\emph{Proof.} The antipode \(w \mapsto -w\) reverses the root order on
\(W\lambda \cup W\lambda^*\) and \(\sigma\) preserves it, so
\(w_0 = -\sigma\) is an anti-automorphism of the weight lattice, hence
of \(J(P)\). Rowmotion sends an ideal \(I\) to the ideal generated by
the minimal elements of \(P \setminus I\); its inverse sends \(I\) to
the ideal whose complement is the filter generated by the maximal
elements of \(I\). An anti-automorphism of \(J(P)\) exchanges ideals
with filters and minimal with maximal elements, so it exchanges the two
recipes. \(\square\)

Checked on all 42 boards \textbf{(C)}: \(w_0 \in W\),
\(w_0(\lambda) = w_0\lambda\), \(w_0Rw_0 = R^{-1}\).

\textbf{Corollary 5.2.} \emph{(i) The Weyl elements reversing \(R\) form
the coset \(w_0\,C_W(R)\) (if \(g, g'\) both reverse \(R\) then
\(g^{-1}g'\) commutes with it); on the rich boards the reverser is
unique. (ii) \(w_0\) commutes with \(R^k\) iff \(R^{2k} = 1\); for \(h\)
even, \(w_0 \in I_{h/2}\) always. (iii) On a self-dual board the
antipode \(\iota\) commutes with \(R^k\) iff \(R^{2k} = 1\), and \(R^k\)
preserves the pair \(\{w,-w\}\) iff \(R^{2k}w = w\); so the antipodal
pairs preserved at lag \(k\) number \(\tfrac12 \mathrm{fix}(R^{2k})\),
and \(R^k\) induces a map on \(W\lambda/\pm\) only when \(R^{2k} = 1\).
On \(E_7\): one pair for \(k = 1, \dots, 8\) and all 28 at \(k = 9\).}

\textbf{Proposition 5.3 (what the reverser does to \(W\)) (C).}
\emph{\(w_0 s_i w_0 = s_{\sigma(i)}\) on every board. The reverser is
central in \(W\), and equal to the antipode, on exactly seven boards of
rank \(\le 7\) --- \(D_4\,\omega_1/\omega_3/\omega_4\),
\(D_6\,\omega_1/\omega_5/\omega_6\), \(E_7\,\omega_7\) --- and on every
other board it is non-central and acts on \(W\) by the diagram flip
(\(\Lambda^k \leftrightarrow \Lambda^{n+1-k}\) on \(A_n\), the two
half-spins on \(D_{\mathrm{odd}}\), \(27 \leftrightarrow \overline{27}\)
on \(E_6\)). On the five self-dual boards with \(-1 \notin W\)
(\(A_3\omega_2\), \(A_5\omega_3\), \(A_7\omega_4\), \(D_5\omega_1\),
\(D_7\omega_1\)) the antipode is a permutation of the weights that
reverses \(R\) and is not a Weyl element, and the isometry group of the
weight configuration is \(W \times \langle -1 \rangle\); on all other
boards it is \(W\).}

\subsection{5.2 The centralizer law and its
exceptions}\label{the-centralizer-law-and-its-exceptions}

\textbf{Proposition 5.4 (= Theorem E(iii)) (C).} \emph{For the 42 boards
of Table 1 with \(W\) enumerated in full: on every board on which no
power of \(R\) is a Weyl element --- the middle representations
\(\omega_k\) (\(1 < k < n\)) of \(A_4, \dots, A_7\), the half-spins of
\(D_5\), \(D_6\) and \(D_7\), \(E_6\) and \(E_7\): 23 boards ---
\(I_k = C_W(R^k)\) as sets for every \(k\), with one exception, the two
half-spins of \(D_5\), where \(|I_3| = |I_5| = 2\) against
\(|C_3| = |C_5| = 1\) and \(|I_4| = 8\) against \(|C_4| = 4\). On the
chains \(I_k = W\). On the vector boards \(D_n\omega_1\), \(|I_k|\)
follows \(\gcd(k, n-1)\) (on \(D_7\): \(12, 72, 384, |W|\) for
\(\gcd = 1, 2, 3, 6\), against centralizers \(2, 12, 16, 384\)).}

The chains and the vectors are explained by Theorems 3.2--3.3 and by
§5.3; the \(D_5\) exception is Theorem F (§6).

\textbf{Remark 5.5.} The statement ``\(I_k = C_k\) on every board'' is
false; the statement ``the unique non-trivial half-turn survivor on a
rich board with \(h\) even is \(-1\) when \(-1 \in W\) and nothing
otherwise'' is also false --- the survivor is \(w_0\) (Corollary
5.2(ii)), which is \(-1\) only when \(-1 \in W\). On the \(D_6\)
half-spins four elements survive the half-turn (\(w_0 = -1\) among them)
and on the \(D_5\) half-spin eight.

\subsection{5.3 The vector boards: the centralizer of the doubled
clock}\label{the-vector-boards-the-centralizer-of-the-doubled-clock}

\textbf{Lemma 5.6.} \emph{If an involution \(\iota\) of the weights
centralizes \(W\) and reverses \(R\), then \(I_k \subseteq C_W(R^{2k})\)
for every \(k\).}

\emph{Proof.} For \(g \in I_k\) let \(h = R^{-k}gR^k \in W\). Then
\(\iota h \iota = R^{k} g R^{-k}\) (since \(\iota R \iota = R^{-1}\) and
\(\iota g \iota = g\)), while \(\iota h \iota = h\) because \(h \in W\)
and \(\iota\) centralizes \(W\). Hence \(R^kgR^{-k} = R^{-k}gR^k\),
i.e.~\(g\) commutes with \(R^{2k}\). \(\square\)

The hypothesis holds on the seven boards with \(-1 \in W\) and on the
five self-dual boards of Proposition 5.3 whose antipode is not a Weyl
element (there \(\iota = -1\) centralizes \(W\) as a linear map); the
inclusion was checked at every lag, with \(W\) enumerated, on the eleven
of these boards other than \(E_7\) and on the \(B_4\) spinor
\textbf{(C)}; on \(E_7\) it follows from Theorem 4.3.

\textbf{Theorem 5.7 (= Theorem E(ii)).} \emph{On the vector boards
\(D_n\,\omega_1\) (\(n \ge 3\), weights \(\pm e_i\)),
\(I_k = C_W(R^{2k})\) for every \(k\).}

\emph{Proof.} \(\subseteq\) is Lemma 5.6 with \(\iota = -1\) (a Weyl
element for \(n\) even, the antipode otherwise; in both cases it
centralizes \(W\) and reverses \(R\) by Theorem 5.1 and Proposition
5.3). For \(\supseteq\) let \(g \in W\) commute with \(R^{2k}\) and put
\(h = R^{-k}gR^k\). Then
\(\iota h \iota = R^{k}gR^{-k} = R^{-k}(R^{2k}gR^{-2k})R^{k} = h\), so
\(h\) is a permutation of \(\{\pm e_i\}\) commuting with \(-1\), i.e.~a
signed permutation, an element of \(W(B_n)\) (the centralizer of \(-1\)
in \(\mathrm{Sym}(2n)\)). Since \(h\) is conjugate to \(g\) in
\(\mathrm{Sym}(2n)\) it has the same sign, and the sign character of
\(\mathrm{Sym}(2n)\) restricted to \(W(B_n)\) is
\((-1)^{\#\text{negative cycles}}\), whose kernel is \(W(D_n)\); so
\(h \in W(D_n) = W\) and \(g \in I_k\). \(\square\)

\textbf{Proposition 5.8 (the transports on the vector boards) (C, with
proof of the first sentence).} \emph{A transport
\(g \mapsto R^{-k}gR^k\) preserves the unsigned cycle type and the sign
of a signed permutation, so it can change the signed cycle type only by
trading two negative \(j\)-cycles for one positive \(2j\)-cycle; no
element of \(W(B_n)\) realises such a trade by conjugation, and \(g\) is
\(W\)-inner (its transport equals \(w^{-1}gw\) for some \(w \in W\)) iff
the signed cycle types of \(g\) and its transport agree. At every lag
coprime to \(n-1\) the shared grammar \(I_k\) is dihedral of order
\(2(n-1)\) and the transport preserves its cyclic half; for \(n\) odd it
is the outer automorphism swapping the two reflection classes (inner in
no group below \(\mathrm{Sym}(2n)\)), and for \(n\) even it is
conjugation by \(R^{h/2}\) or by \(w_0R^{h/2}\) --- because on
\(C_W(R^{2k})\) the transport equals conjugation by any odd power of
\(R^k\), and \(h/2 = n-1\) is such a power exactly when \(n\) is even.
The non-inner counts: \(D_3\), 2 of 4; \(D_5\), 4 of 8 at odd lags and
18 of 32 at lags 2, 6; \(D_7\), 6 of 12 at lags coprime to 6 and 252 of
384 at lags 3, 9; none on \(D_4\), \(D_6\).}

\textbf{Remark 5.9 (a false lemma).} The identity
\(I_k = \bigcup_{w \in W} C_W(R^kw^{-1})\) is tempting and false. The
inclusion \(\supseteq\) holds; equality holds iff every transport is
\(W\)-conjugate to its element. It fails on the \(B_3\) spinor (8 of the
16 elements of \(I_3\)), on \(A_3\omega_2\) and on the odd \(D_n\)
vectors --- exactly the non-inner transports of Proposition 5.8. Across
the 98 (board, lag) pairs of the 42 boards on which
\(I_k \ne C_W(R^k)\), on 83 a single Weyl element realises the transport
and on 15 some transports leave their \(W\)-class.

\section{6. The 4-cube}\label{the-4-cube}

\subsection{6.1 Locating the exception}\label{locating-the-exception}

The \(D_5\) half-spin (16 weights, \(h = 8\), \(R\) two free 8-cycles)
shares its poset and its rowmotion with the \(B_4\) spinor board --- the
16 vertices of the 4-cube under \(W(B_4) = \mathbb{F}_2^4 \rtimes S_4\)
of order 384 --- and \(W(B_4) < W(D_5)\) \textbf{(C)}. The half-turn
anomaly is already present on \(B_4\): \(|I_4| = 8\) against
\(|C_4| = 4\), with the same four extra elements --- two signed 4-cycles
\(e_1 \to -e_4 \to -e_2 \to e_3 \to e_1\) and its inverse, and two
involutions pairing \((e_1, e_4)\) with \((e_2, e_3)\). Only the lag-3
survivor needs the fifth axis: it is the coordinate permutation
\((e_1e_3)(e_2e_5)\), the one survivor that mixes the clock's two
orbits. What the anomaly is \emph{not} \textbf{(C)}: not a hidden
\(\mathbb{F}_2\)-affinity (no power of the \(D_5\), \(D_6\) or \(D_7\)
clock is an affine map of the even code; \(D_4\)'s half-turn is, being
the linear half-turn of Theorem 3.3); not a near-miss of linearity (the
altitude of \(R^4\) against \(W(D_5)\) is 8 of 16, of \(R^3\) 7); not an
overgroup effect (\(\langle W(D_5), R^4 \rangle = A_{16}\); contrast the
\(B_3\) spinor, whose half-turn signature --- shared grammar 16 against
centralizer 8 --- comes from \(W(D_4) \supset W(B_3)\) containing the
\(D_4\) half-turn); and not specific to the 4-cube in the naive sense
(the \(B_5\) spinor is clean, \(I_k = C_k\) for every \(k\)).

\subsection{6.2 The mechanism}\label{the-mechanism}

\textbf{Lemma 6.1.} \emph{Let \(w \in W\) and suppose
\(Y := \mathrm{Fix}(R^kw^{-1})\) has trivial pointwise stabilizer in
\(W\). Then \(I_k \cap \mathrm{Stab}_W(Y) = C_W(R^kw^{-1})\), and on
this set the transport is conjugation by \(w\).}

\emph{Proof.} Put \(c = R^kw^{-1}\), so \(R^k = cw\) and \(c\) is the
identity on \(Y\). If \(g \in W\) commutes with \(c\) then
\(R^{-k}gR^k = w^{-1}c^{-1}gcw = w^{-1}gw \in W\), so \(g \in I_k\), and
\(g\) preserves \(Y = \mathrm{Fix}(c)\). Conversely let \(g \in I_k\)
preserve \(Y\) and let \(h = R^{-k}gR^k \in W\); then
\(c^{-1}gc = whw^{-1} \in W\). For \(y \in Y\), \(c(y) = y\) and
\(g(y) \in Y\), so \(c^{-1}gc(y) = c^{-1}g(y) = g(y)\): the two Weyl
elements \(c^{-1}gc\) and \(g\) agree on \(Y\), hence are equal by the
hypothesis on \(Y\), i.e.~\(g\) commutes with \(c\); and then
\(h = w^{-1}gw\). \(\square\)

\textbf{Theorem 6.2 (= Theorem F).} \emph{On the \(B_4\) spinor board,
\(R^4\) agrees with \(-\tau\), \(\tau = (e_2e_3)(-e_4)\), on all of
\(O_2\) and with no Weyl element on \(O_1 \ni \lambda\). Consequently
\(I_4 = C_W(c)\) for the correction \(c = (-\tau)R^4\), an involution
that is the identity on \(O_2\) and fixed-point-free on \(O_1\);
\(|I_4| = 8\), the four extras being the elements commuting with \(c\)
but not with \(\tau\); the transport on \(I_4\) is conjugation by
\(\tau\); and \([g, R^4] = -1\) for each extra \(g\).}

\emph{Proof.} Grade the weights by rank \(= \mathrm{wt}(x)\). For a Weyl
element \(w = (\pi, v)\) put \(u := \pi^{-1}(v)\); since a coordinate
permutation preserves Hamming weight,

\[\mathrm{rank}(w(x)) = \mathrm{wt}(\pi(x) \oplus v) = \mathrm{wt}(x \oplus u), \qquad \Delta_w(x) := \mathrm{rank}(w(x)) - \mathrm{rank}(x) = \mathrm{wt}(u) - 2\,|x \cap u| .\]

The orbit \(O_1\) of \(\lambda = 0\) has ranks
\(0, 4, 3, 3, 2, 2, 1, 1\) in the order
\(\lambda, R\lambda, R^2\lambda, \dots\) (the \emph{full-height} orbit:
\(R\) sends \(\lambda\) straight to \(w_0\lambda\) by Lemma 3.1(a), and
the ranks then descend in pairs), and \(R^4\) pairs clock positions
\(j\) and \(j+4\), whose ranks differ by exactly 2. If
\(R^4|_{O_1} = w|_{O_1}\) then \(\Delta_w(x) = \pm 2\) on \(O_1\); at
\(x = 0\), which \(R^4\) moves up to rank 2,
\(\Delta_w(0) = \mathrm{wt}(u) = 2\); then
\(\Delta_w(x) = 2 - 2|x \cap u| \in \{2, 0, -2\}\) and
\(|\Delta_w| = 2\) forces \(|x \cap u| \ne 1\),
i.e.~\(\langle x, u \rangle = 0\), for every \(x \in O_1\). But \(O_1\)
contains \(0\) and spans \(\mathbb{F}_2^4\), so \(u = 0\), contradicting
\(\mathrm{wt}(u) = 2\). On \(O_2\) the ranks stay in \(\{1,2,3\}\) and
\(R^4\) shifts each by \(\pm 1\); \(u = 1110\) gives
\(\Delta = 3 - 2|x \cap u| \in \{1,-1\}\) for all \(x \in O_2\), and the
Weyl element with this \(u\) and the matching translation is \(-\tau\),
which therefore agrees with \(R^4\) on \(O_2\) (uniquely, \(O_2\)
spanning). \(O_2\) affinely spans \(\mathbb{F}_2^4\), so its pointwise
stabilizer in \(W\) is trivial, and Lemma 6.1 with \(Y = O_2\),
\(w = -\tau\) gives \(I_4 \cap \mathrm{Stab}_W(O_2) = C_W(c)\). Finally
the sixteen Weyl elements that swap the two orbits (of the 384, sixteen
stabilise each orbit, sixteen swap them, and the rest mix them) lie in
no \(I_k\) (checked \textbf{(C)}; by hand, \(\mathrm{Stab}_W(O_1)\) is
sixteen affine maps whose image in \(S_4\) is the dihedral stabilizer of
the pairing \(\{14|23\}\), kernel \(\pm 1\), and \(I_4\) is the preimage
of its Klein four-group), so \(I_4 = C_W(c)\), of order 8. The remaining
assertions are Lemma 6.1's transport statement and direct computation.
\(\square\)

The lag-3 survivor of \(D_5\) falls to the same lemma with the nearest
Weyl element \(w_3\) (unique at Hamming distance 7; its agreement set of
nine points determines \(W(D_5)\)): \((e_1e_3)(e_2e_5)\) commutes with
\(R^3w_3^{-1}\), and its lag-5 partner is its transport \textbf{(C)}.
The anomaly sits on \(B_4\), where the reverser \(-1\) is central, so it
is not tied to \(D_5\)'s non-central reverser.

\textbf{Lemma 6.3 (dimension-free form).} \emph{On the \(B_n\) spinor
board, if an \(R\)-orbit \(O\) contains \(\lambda = 0\), spans
\(\mathbb{F}_2^n\), and \(R^{h/2}\) shifts every rank on \(O\) by a
constant magnitude \(m \ge 2\), then \(R^{h/2}|_O\) is not the
restriction of a Weyl element.} (\emph{Proof.} If \(w = (\pi, v)\)
agreed with \(R^{h/2}\) on \(O\), then at \(\lambda = 0\) the shift is
\(\mathrm{wt}(u) = m\), and \(|\Delta_w(x)| = m\) forces
\(|x \cap u| \in \{0, m\}\): every \(x \in O\) misses \(u\) or contains
it. These two parallel affine subspaces have codimension \(m\), so their
union spans a subspace of dimension \(n - m + 1 < n\), contradicting
that \(O\) spans \(\mathbb{F}_2^n\). \(\square\)) On the full-height
orbit of \(B_n\) the half-turn shifts rank by exactly \(\pm n/2\) for
every even \(n\) (Proposition 6.13); the clock acts freely on the
\(2^n\) weights with orbits of size \(2n\) only when \(2n \mid 2^n\),
i.e.~\(n\) a power of 2, so the clean two-orbit dichotomy of \(B_4\)
next arises, if at all, on \(B_8\) (256 weights, 16 orbits of 16). It
does not (Proposition 6.7, §6.4).

\subsection{6.3 Iterated shared
grammars}\label{iterated-shared-grammars}

For a board and a lag \(k\) define \(P_1(k) = W\) and
\(P_{j+1}(k) = \{g \in P_j : R^{-k}gR^k \in P_j\}\), so \(P_2 = I_k\);
let \(d(k)\) be the first \(j\) with \(P_{j+1} = P_j\) and
\(P_\infty(k)\) the stable value. (The definition mimics the Clifford
hierarchy of quantum computation, with \(W\) in the role of the Clifford
group; we use it only as a chain of subgroups.)

\textbf{Lemma 6.4.}
\emph{\(P_j(k) = \bigcap_{0 \le i < j} R^{ik} W R^{-ik}\). Hence
\(P_\infty(k)\) is the largest subgroup of \(W\) normalised by \(R^k\),
and it contains the core of \(W\) in \(\langle W, R^k \rangle\).}

\textbf{Lemma 6.5.} \emph{\(P_3(k) = P_2(k)\) iff \(I_k = I_{-k}\) iff
some (equivalently every) reverser of \(R\) in \(W\) normalises
\(I_k\).}

\emph{Proof.} \(P_3 = W \cap R^kWR^{-k} \cap R^{2k}WR^{-2k}\), so
\(P_3 = P_2\) iff \(I_k \subseteq R^{2k}WR^{-2k}\), iff (conjugating by
\(R^{-k}\)) \(I_{-k} \subseteq I_k\), iff \(I_{-k} = I_k\) since the two
have equal order. For a reverser \(r\),
\(rI_kr = W \cap rR^kWR^{-k}r = W \cap R^{-k}WR^{k} = I_{-k}\).
\(\square\)

\textbf{Proposition 6.6 (C, 44 boards, 278 (board, lag) pairs).}
\emph{\(d(k) = 1\) exactly where \(I_k = W\) (the chains; the \(D_n\)
vectors at the half-turn). \(d(k) = 2\) everywhere else --- two instants
already share what all instants share --- with exactly one exception,
the \(D_5\) half-spin at lags 3 and 5, where the survivor is carried to
a different element, no reverser normalises \(I_3\), and the chain drops
once more: \(1920 \supset 2 \supset 1\). The stable groups: \(C_W(R^k)\)
on the rich boards (trivial off the half-turn; at the half-turn
\(\{\pm 1\}\) where \(-1 \in W\), except on the \(D_6\) half-spins,
where it has order 4, Remark 5.5); \(C_W(R^{2k})\) on the \(D_n\)
vectors; the dihedral \(I_4\) of order 8 on the 4-cube and the \(D_5\)
half-spin at the half-turn; on \(E_7\), trivial at every lag but the
half-turn, where it is \(\{\pm 1\}\). Wherever \(P_\infty\) is neither
\(W\) nor of order \(\le 2\) (44 cases) it is not normal in \(W\); where
\(\langle W, R \rangle\) could be enumerated (\(D_3\omega_1\),
\(D_4\omega_1\), the \(B_3\) spinor: \(A_6\), \(A_8\), \(A_8\)) the core
of \(W\) in \(\langle W,R\rangle\) is trivial, below \(P_\infty(1)\) (of
orders 4, 6 and 1; equality in the last case).}

\subsection{6.4 The 8-cube, the 16-cube, and the commutation
theorem}\label{the-8-cube-the-16-cube-and-the-commutation-theorem}

The clock acts freely on the \(2^n\) weights of the \(B_n\) spinor
board, with orbits of size \(h = 2n\), only when \(2n \mid 2^n\),
i.e.~when \(n\) is a power of 2; \(B_8\) is therefore the next board on
which the two-orbit mechanism of Theorem 6.2 could recur, and the last
one within reach (\(B_{16}\) has 65,536 weights).
\(W(B_8) = \mathbb{F}_2^8 \rtimes S_8\) and \(W(D_9)\) were enumerated
in the affine form \(x \mapsto \pi(x) \oplus v\) on the 256 sign
patterns (for \(D_9\), the even code with its parity bit dropped), never
as permutation tuples; membership of a transport \(R^{-k}gR^k\) in \(W\)
is decided by its values at \(0\) and the unit vectors and then
confirmed on all 256 points. \emph{Notation.} For a product \(\pi\) of
disjoint transpositions of the coordinates \(e_2, \dots, e_{n-1}\) we
write \(-\pi\) for the signed permutation
\(x \mapsto \pi(x) \oplus (e_1 + \dots + e_{n-1})\): it negates
\(e_1, \dots, e_{n-1}\), permutes them by \(\pi\), and fixes \(e_n\); on
\(B_4\), \(-(e_2e_3)\) is the element \(-\tau\) of Theorem 6.2. The
correction is written \(c = w^{-1}R^{h/2}\) rather than Lemma 6.1's
\(R^{h/2}w^{-1}\); since
\(\mathrm{Fix}(R^{h/2}w^{-1}) = w\,\mathrm{Fix}(w^{-1}R^{h/2})\) and in
every case below \(w\) preserves its agreement set
\(Y = \mathrm{Fix}(c)\) (the stabilizers computed are
\(\langle -1, w \rangle\), or \(w\) commutes with \(R^{h/2}\)), the two
fixed sets coincide, Lemma 6.1 applies to \(Y\) verbatim, and
\(C_W(R^{h/2}w^{-1}) = w\,C_W(c)\,w^{-1}\) has the order stated.

\textbf{Proposition 6.7 (C).} \emph{On the \(B_8\) spinor board (\(256\)
weights, \(|W| = 2^8 \cdot 8! = 10{,}321{,}920\), \(h = 16\)):}

\emph{(i) \(R\) has sixteen free orbits of sixteen. The orbit
\(O_1 \ni \lambda\) is the full-height orbit, with ranks
\(0, 8, 7, 7, \dots, 1, 1\) in the order
\(\lambda, R\lambda, R^2\lambda, \dots\), and \(R^8\) shifts every rank
on it by exactly \(\pm 4\); by Lemma 6.3 no Weyl element agrees with
\(R^8\) on \(O_1\).}

\emph{(ii) Exactly two orbits carry the half-turn as a Weyl element,
both by the same element \(w = -(e_2e_3)(e_4e_5)(e_6e_7)\), which
negates \(e_1, \dots, e_7\) and fixes \(e_8\) --- the analogue of
\(-\tau\). It is the unique Weyl element nearest to \(R^8\), agreeing
with it on \(52\) of the \(256\) weights. Each carrying orbit is
preserved by \(-1\). Every orbit affinely spans \(\mathbb{F}_2^8\).}

\emph{(iii) Lemma 6.1 applies with \(Y = \mathrm{Fix}(w^{-1}R^8)\)
(\(52\) points, spanning), but the correction \(c = w^{-1}R^8\) is not
an involution (cycle type \(8^8\,4^{12}\,3^{12}\,2^{28}\,1^{52}\)) and
\(C_W(c) = \{\pm 1\}\). Hence \(I_8 = C_8 = \{\pm 1\}\), and
\(I_k = C_k\) for every \(k = 1, \dots, 15\). The same holds at every
lag on the \(D_9\) half-spin (the same clock;
\(|W(D_9)| = 92{,}897{,}280\)) and on the \(B_6\) and \(B_7\) spinor
boards.}

\emph{Consequently, among the spinor boards \(B_n\) with
\(3 \le n \le 8\), the shared grammar exceeds the centralizer only on
\(B_4\) and, through the overgroup \(W(D_4)\), on \(B_3\) (§6.1). The
mechanism of Theorem 6.2 --- the half-turn agreeing with a Weyl element
on an orbit that avoids the extreme ranks --- recurs on \(B_8\); its
consequence does not, because the correction's centralizer collapses to
\(\{\pm 1\}\).}

\textbf{Remark 6.8.} On \(B_8\) the two carrying orbits are exactly the
orbits whose ranks stay in \(\{3,4,5\}\), equivalently those on which
every half-turn rank-shift is \(\pm 1\); on \(B_4\) the carrying orbit
\(O_2\) likewise has ranks \(\{1,2,3\}\). On \(B_{16}\) the sixteen
carrying orbits are again exactly the band \(\{7,8,9\}\) (Proposition
6.11). This is an observation on three boards; the orbit of \(\lambda\)
itself is described exactly in Proposition 6.13. The expectations that
the phenomenon of Theorem 6.2 recurs on \(B_8\), that it marks the
boards with a free clock (``\(n\) a power of 2''), that \(-1\) exchanges
the two carrying orbits, and that \(D_9\) carries a survivor at a lag
coprime to \(16\) were all inverted by the computation.

\textbf{Lemma 6.9 (the half-turn criterion).} \emph{Let \(m = h/2\),
\(w \in W\) and \(c = w^{-1}R^m\). (i) \(c^2 = 1\) if and only if
\(R^m w R^m = w^{-1}\); for an involution \(w\) this says
\(w \in C_W(R^m)\). (ii) If \(w \in C_W(R^m)\) then \(w\), \(c\) and
\(R^m\) pairwise commute; for \(g \in W\) any two of ``\(g\) commutes
with \(w\)'', ``with \(R^m\)'', ``with \(c\)'' imply the third; and when
\(\mathrm{Fix}(c)\) has trivial pointwise stabilizer,
\(I_m \cap \mathrm{Stab}_W(\mathrm{Fix}\,c) = C_W(c)\) with
\(\big(I_m \cap \mathrm{Stab}_W(\mathrm{Fix}\,c)\big) \setminus C_W(R^m) = C_W(c) \setminus C_W(w)\).}

\emph{Proof.} (i)
\(c^2 = w^{-1}R^m w^{-1}R^m = 1 \iff R^m w^{-1} R^m = w \iff R^m w R^m = w^{-1}\),
since \(R^m\) is an involution. (ii)
\(cw = w^{-1}R^m w = w^{-1} w R^m = R^m = wc\), and both commute with
\(R^m = wc\); if \(g\) commutes with two of \(w\), \(c\), \(wc\) it
commutes with the third. The rest is Lemma 6.1 together with the
observation that for \(g \in C_W(c)\), \(g \notin C_W(R^m)\) if and only
if \(g \notin C_W(w)\). \(\square\)

\textbf{Proposition 6.10 (C).} \emph{On the spinor boards
\(B_3, \dots, B_8\), \(I_{h/2} \supsetneq C_{h/2}\) if and only if a
Weyl element nearest to \(R^{h/2}\) (one maximising the agreement set
\(\mathrm{Fix}(w^{-1}R^{h/2})\)) commutes with \(R^{h/2}\). On \(B_4\)
the nearest element is \(-\tau\), it commutes with \(R^4\), \(c\) is an
involution and the four extras are \(C_W(c) \setminus C_W(-\tau)\). On
\(B_5, \dots, B_8\) the nearest element is unique, of the shape
\(-(e_2e_3)(e_4e_5)\cdots\) (with \(e_n\) fixed for \(n\) even, and
\(e_{n-1}, e_n\) negated for \(n\) odd), does not commute with
\(R^{h/2}\), \(c\) is not an involution, and \(C_W(c) = \{\pm 1\}\); on
\(B_8\) the stabilizer of the 52-point agreement set is
\(\langle -1, w \rangle\). On \(B_3\) the four nearest elements all
commute with \(R^3\) and \(|I_3| = 16 > |C_3| = 8\), so the equivalence
holds there as well; but no element of \(W(B_3)\) agrees with \(R^3\) on
an affinely spanning set, Lemma 6.1 does not apply, and the excess is
explained instead by the overgroup \(W(D_4) \ni R^3\) (§6.1).}

\textbf{Proposition 6.11 (C, the 16-cube at the orbit level).} \emph{On
the \(B_{16}\) spinor board (\(65{,}536\) weights, \(h = 32\); \(W\) is
not enumerated): \(R\) has \(2{,}048\) free orbits of \(32\); the orbit
of \(\lambda\) is the full-height orbit and \(R^{16}\) shifts every rank
on it by exactly \(\pm 8\); exactly sixteen orbits carry the half-turn
as a Weyl element, all by the same
\(w = -(e_2e_3)(e_4e_5)\cdots(e_{14}e_{15})\) with \(e_{16}\) fixed, and
they are exactly the orbits whose ranks lie in \(\{7,8,9\}\); \(w\) does
not commute with \(R^{16}\); \(C_W(R^{16}) = \{\pm 1\}\); the agreement
set \(Y = \mathrm{Fix}(w^{-1}R^{16})\) has \(1{,}900\) points and spans,
\(\mathrm{Stab}_W(Y) = \langle -1, w \rangle\) and
\(C_W(w^{-1}R^{16}) = \{\pm 1\}\), so by Lemma 6.1
\(I_{16} \cap \mathrm{Stab}_W(Y) = \{\pm 1\}\). Whether
\(I_{16} = \{\pm 1\}\) is not decided.} (Rowmotion was computed
combinatorially on the shifted staircase and validated against the
\(B_4\) and \(B_8\) clocks; the Weyl element agreeing with \(R^{16}\) on
an orbit was found by matching columns of the difference vectors, the
stabilizer and the centralizers by backtracking over coordinate
permutations with every solution verified on all \(65{,}536\) weights.
\(1{,}987\) of the \(2{,}048\) orbits affinely span
\(\mathbb{F}_2^{16}\), every carrying orbit among them.)

\textbf{Lemma 6.12 (rowmotion in excess coordinates).} \emph{Write a
weight of the \(B_n\) spinor board by its plus positions
\(p_1 < \dots < p_r\) and put \(d_k = p_k - k\), so
\(0 \le d_1 \le \dots \le d_r \le n - r\) (the empty sequence is
\(w_0\lambda\)). Call \(k\) a gap row if \(d_k > d_{k-1}\)
(\(d_0 := 0\)). Then the excess sequence of \(R(x)\) is}

\[d'_{k'} = \min\Big(\{\,d_k - 1 : k \text{ a gap row},\ k \ge k'\,\} \cup \{\,n - r - 1 : \text{if } d_r < n - r \text{ and } k' \le r + 1\,\}\Big),\]

\emph{for every \(k'\) for which the set is non-empty; rows with an
empty set are absent.}

\emph{Proof.} Identify \(P_\lambda\) with the shifted staircase
\(\delta_n\) of §2: cells \((k, c)\) with \(1 \le k \le c \le n\) and
colour \(c\), \((k, c)\) covering \((k, c+1)\) and \((k-1, c-1)\) ---
the \(k\)-th plus sign, counted from the left, having reached position
\(c\) --- so that the ideal of the weight with plus positions
\(p_1 < \dots < p_r\) is \(\{(k, c) : k \le r,\ c \ge p_k\}\), a cell
\((k, c) \notin I\) with \(k \le r\) has \(c < p_k\) and is minimal in
\(P \setminus I\) iff its lower covers lie in \(I\), i.e.~iff
\(c = p_k - 1\) and (\(k = 1\) or \(p_k - 1 \ge p_{k-1} + 1\)), which is
\(d_k > d_{k-1}\); the only candidate in row \(r+1\) is \((r+1, n)\),
minimal iff \(r = 0\) or \(p_r \le n - 1\), i.e.~\(d_r < n - r\); rows
\(\ge r + 2\) contain no minimal element. The down-closure of \((k, c)\)
is \(\{(k', c') : k' \le k,\ c' \ge c - k + k'\}\), so the threshold of
\(R(I)\) in row \(k'\) is the minimum of \(c - k + k'\) over the
generators \((k, c)\) with \(k \ge k'\); subtracting \(k'\) gives the
formula. \(\square\)

\textbf{Proposition 6.13 (the orbit of \(\lambda\)).} \emph{For
\(1 \le j \le \lfloor n/2 \rfloor\),
\(R^{2j}\lambda = [n] \setminus \{n, n-2, \dots, n-2j+2\}\), and for
\(0 \le j \le \lfloor (n-1)/2 \rfloor\),
\(R^{2j+1}\lambda = [n] \setminus \{n-1, n-3, \dots, n-2j+1\}\) (weights
written as sets of minus positions). Hence the ranks around the orbit of
\(\lambda\) read \(0, n, n-1, n-1, \dots, 1, 1\), and for even \(n\) the
half-turn \(R^n\) shifts every rank on this orbit by exactly
\(\pm n/2\).}

\emph{Proof.} In excess coordinates \(R^{2j}\lambda\) is
\(d_k = n - 2j + k\) (\(k \le j\)) and \(R^{2j+1}\lambda\) is
\(d_k = n - 1 - 2j + k\); both are strictly increasing, so every row is
a gap row, and one checks from Lemma 6.12 that the first has no cap and
loses \(1\) in every entry, while the second has the cap \(n - j - 1\),
which is larger than every \(d_k - 1\) and is appended as a new row. The
base \(R\lambda = w_0\lambda\) is the empty sequence. Ranks follow by
counting, and the second half of the orbit from
\(R^{-k}\lambda = -R^{k+1}\lambda\) (Theorem 5.1 with
\(w_0\lambda = R\lambda\), Lemma 3.1). \(\square\)

\textbf{Theorem 6.14 (why four; = Theorem G).} \emph{Let \(n \ge 4\) be
even and \(w_n = -(e_2e_3)(e_4e_5)\cdots(e_{n-2}e_{n-1})\), the signed
permutation negating \(e_1, \dots, e_{n-1}\) and fixing \(e_n\). Then
\(w_n\) commutes with the half-turn \(R^n\) if and only if \(n = 4\).}

\emph{Proof.} Write \(O_1\) for the orbit of \(\lambda\). (a) \(w_n\)
preserves \(O_1\): by Proposition 6.13 and the identity
\(R^{-k}\lambda = -R^{k+1}\lambda\), a direct computation gives
\(w_n(R^{2j}\lambda) = R^{2-2j}\lambda\) and
\(w_n(R^{2j+1}\lambda) = R^{-(2j+1)}\lambda\); the maps
\(i \mapsto 2 - i\) and \(i \mapsto -i\) on \(\mathbb{Z}/2n\) commute
with \(i \mapsto i + n\) when \(n\) is even, so \(w_n R^n = R^n w_n\) on
\(O_1\). (b) If \(R^n\) agrees with the involution \(w_n\) on an orbit
\(O\), then \(w_nR^nx = x = R^nw_nx\) for \(x \in O\). (c) For \(n = 4\)
the board is \(O_1 \cup O_2\) and \(R^4 = -\tau = w_4\) on \(O_2\)
(Theorem 6.2), so \(w_4 \in C_W(R^4)\). (d) For even \(n \ge 6\) let
\(e_2 = \{2\}\). Iterating Lemma 6.12 along the orbit of \(e_2\) gives,
for \(2 \le j \le n/2 - 1\), the excess sequences
\((n-1-2j, n-2j+2, \dots, n-j)\) at \(R^{2j}e_2\) and
\((n-2-2j, n-2j+1, \dots, n-j-1)\) at \(R^{2j+1}e_2\), and then
\(R^n(e_2) = \{1, 2\} \cup \{5, 7, \dots, n-1\}\); along the orbit of
\(\{3, n\}\) (whose first step is \(e_4\)) it gives
\(R^n(\{3, n\}) = [n] \setminus \{4, 6, 7, 9, 11, \dots, n-1\}\) for
\(n \ge 8\) and \(\{1, 2, 3, 5, 6\}\) for \(n = 6\). Since
\(\pi = (2\,3)(4\,5)\cdots\) swaps \(2\) and \(3\) and
\(3 \notin R^n(e_2)\), we get
\(2 \in w_n(R^n e_2) = \pi(R^ne_2)\,\triangle\,[n-1]\); and
\(w_n(e_2) = [n-1] \setminus \{3\} = -\{3, n\}\), so
\(R^n(w_n e_2) = -R^n(\{3, n\}) \not\ni 2\) because
\(2 \in R^n(\{3, n\})\). Hence \(w_nR^ne_2 \ne R^nw_ne_2\). \(\square\)

The theorem is checked on the actual clocks for even \(n \le 16\), Lemma
6.12 on every weight for \(n \le 12\), and every excess sequence named
in the proof for \(6 \le n \le 16\) \textbf{(C)}. Together with Lemma
6.9 it is the reason for the \(4\)-cube: the mechanism of Theorem 6.2 is
present on every free clock, and its consequence requires a commutation
that only the smallest board allows, because there the two kinds of
orbit on which commutation is forced exhaust the board.

\section{7. Homomesy of the label
counts}\label{homomesy-of-the-label-counts}

A statistic on weights is \emph{homomesic} under \(R\) if its average is
the same on every \(R\)-orbit {[}PR15{]}. Rush and Wang proved that on
every minuscule poset the ideal cardinality and the antichain
cardinality are homomesic under rowmotion, and that the per-colour ideal
cardinalities are {[}RW15{]}; on the weights the latter says every orbit
has mean weight zero.

A weight's label at node \(i\) is \(+1\) iff its ideal has a maximal
element of colour \(i\), \(-1\) iff the complement has a minimal element
of colour \(i\), and \(0\) otherwise. Hence \(\#(+1) = |\max I|\), the
antichain cardinality; \(\#(-1) = |\min(P \setminus I)| = |\max R(I)|\);
and \(\#(0) = r - \#(+1) - \#(-1)\), \(r\) the rank.

\textbf{Proposition 7.1.} \emph{On every simply-laced minuscule board
the three ternary label counts are homomesic under \(R\), with orbit
means}

\[\#(+1) = \#(-1) = \frac{|P|}{h}, \qquad \#(0) = r - \frac{2|P|}{h} .\]

\emph{The toggle size \(|I \,\triangle\, R(I)|\) is homomesic with mean
\(2|P|/h\).}

\emph{Proof.} The first two counts are the antichain cardinality of
\(I\) and of \(R(I)\), homomesic with mean \(|P|/h\) by {[}RW15{]} (the
equality of the two means is also the 0-mesy of \(T^+ - T^-\) of
{[}DHPP23{]}, valid on every finite poset); the third is a linear
combination. The toggle size is
\(|\max I| + |\min(P \setminus I)| = \mathrm{ac}(I) + \mathrm{ac}(R(I))\),
and a sum of homomesic statistics, one composed with the map, is
homomesic with the sum of the means {[}DHPP23, Thm 3.20 and the closure
properties stated there{]}. \(\square\)

Checked on all 42 boards \textbf{(C)}; \(E_7\): \(3/2, 3/2, 4\);
\(E_6\): \(4/3, 4/3, 10/3\); the \(D_5\) spinor: \(5/4, 5/4, 5/2\).
\textbf{Remark 7.2.} On every non-chain board the number of distinct
colours toggled, the per-weight kept count of Proposition 4.1, and the
number of lattice elements below a weight are \emph{not} homomesic; on
the 56 the ternary counts are not homomesic under \(\mathrm{pr}\), and
are homomesic under the sheet clock \(R_6\) on each 27-sheet but not on
the 56 (the poles are fixed points).

\section{8. The computational model}\label{the-computational-model}

The objects of §§3--7 were first met inside an explicit model, built in
plain Python from generators in one coordinate system, of the following
classical chain. We record what was constructed and how each
identification was made; nothing in this section is claimed as new group
theory.

\textbf{8.1 The roof.} \(W(E_8)\) has order \(696{,}729{,}600\); its
rotation subgroup \(W^+(E_8)\) has index 2 and
\(W^+(E_8)/\{\pm 1\} \cong O_8^+(2)\), simple of order
\(174{,}182{,}400\) {[}Con71, CS99, ATLAS{]}.
\(W(E_7) \cong 2 \times Sp_6(2)\), and \(Sp_6(2)\) acts on the 28
bitangents of a plane quartic {[}Dol12{]}. The 56 weights of the
minuscule representation of \(E_7\) are the roots \(r\) of \(E_8\) with
\(\langle r, \alpha \rangle = 1\) for a fixed root \(\alpha\); the
\(E_7\) root system is \(\alpha^\perp\) {[}Bou02{]}.

\textbf{Proposition 8.1 (C).} \emph{Two families of finite groups ---
the chain \(PSL(2,7) \supset C_6 \times C_2,\ A_5,\ S_4,\ A_4,\ C_2\)
and the tower \(SL(2,7)\), \(2\cdot I = SL(2,5)\), \(GL(2,3)\),
\(2\cdot O\), \(2\cdot T = SL(2,3)\), \(C_3 \times D_4\),
\(2\cdot W(E_6)\), \(2 \times PGL(2,7)\) together with the split covers
\(A_5 \times C_2\), \(A_4 \times C_2\) --- are exhibited by explicit
witnesses as subgroups of \(W(E_8)\) acting on the 240 roots, the first
inside a vector-type subgroup \(C \cong Sp_6(2)\), the second inside
\(H \cong 2\cdot Sp_6(2)\). The group \(K = \langle C, H \rangle\) has
order \(348{,}364{,}800 = |W(E_8)|/2\) and every generator has
determinant \(+1\) (exact Bareiss elimination), so \(K = W^+(E_8)\);
\(K\) is perfect, and the 21 witness generators alone regenerate it. The
centre of \(H\) is \(\{\pm 1\}\): both families hang over a single
central \(C_2\).} (44 checks.)

\textbf{8.2 Triality and the still point.}
\(\mathrm{Out}(O_8^+(2)) \cong S_3\); the fixed group of a triality
automorphism of order 3 is \(G_2(2) \cong U_3(3){:}2\) of order
\(12{,}096\) {[}Con71, Kle87, Wil09, §4.7{]}.

\textbf{Proposition 8.2 (C).} \emph{Let \(\bar C \cong Sp_6(2)\) be the
image of \(C\) in \(\Omega = O_8^+(2)\) (its vector shadow) and
\(\tau'\) the outer automorphism of \(\Omega\) realised in the model.
The intersection \(\bar C \cap \tau'(\bar C)\), computed by sifting all
\(1{,}451{,}520\) elements of \(\bar C\) through a base and strong
generating set of its image, has order \(12{,}096\), derived subgroup of
order \(6{,}048\) and element orders \(\{1,2,3,4,6,7,8,12\}\): it is
\(G_2(2)\). There is an intertwiner \(y\), unique up to a 2-dimensional
solution space, with \(\tau'' = \mathrm{inn}(y^{-1}) \circ \tau'\) of
exact order 3 fixing this \(G_2(2)\) pointwise. Inside it, \(PSL(2,7)\)
occurs as a single class of 36 subgroups, each with normaliser
\(PGL(2,7)\), all of the class transitive on the 28 bitangents; the
still point lifts split to \(2 \times G_2(2)\) in the spin cover.} (19
checks; the split lift contradicted our expectation of a twisted lift.)

\textbf{Proposition 8.3 (C).} \emph{The odd half-spin module \(S^-\)
carries a unique quadratic form \(Q^-\) of plus type; letting
\(\Omega = O_8^+(2)\) act on the \(360 = 3 \times 120\) nonsingular
points of \(V \oplus S^+ \oplus S^-\) and defining \(\Phi\) blockwise
(\(T''^{-1}\) on \(V\), where \(T''\) is the matrix inducing \(\tau''\)
on \(V\); the unique intertwiner \(S^+ \to S^-\), the closing map on
\(S^-\)), \(\Phi\) has exact order 3, conjugates the generators as
\(\tau''\) does, and
\(|\langle \Omega, \Phi \rangle| = 522{,}547{,}200 = 3\,|O_8^+(2)|\):
this is \(O_8^+(2){:}3\) as an explicit permutation group on 360 points.
The still point is the pointwise stabilizer of the \(\Phi\)-orbit
\(\{v, \Phi v, \Phi^2 v\}\) --- one vector, one spinor, one co-spinor.}
(16 checks.)

The composition of \(\Phi\) with a bridge element of order 7 has order
21, and with the image of a Coxeter element of \(W(E_8)\) has order 18:
these are the outer classes \(21A\) and \(18A/B\) of §8.4. The number 18
is the Coxeter number of \(E_7\); its appearance here is a fact about
\(O_8^+(2){:}3\), distinct from Rush--Shi's theorem on the board.

\textbf{8.3 The board under the still point.}

\textbf{Proposition 8.4 (C).} \emph{The 56 weights \(\Lambda\) are the
56 nonsingular vectors \(u\) with \(B(u,v) = 1\) for a fixed nonsingular
\(v\), equivalently the 56 roots of \(E_8\) with
\(\langle r, \alpha \rangle = 1\), matched bijectively by Dynkin labels
with all 84 covering relations of the weight lattice equal to
subtraction of a simple root. Under this identification \(\iota = -1\)
is the transvection \(t_v : u \mapsto u + v\) (Dickson invariant 1,
outside \(\Omega\)), while \(R\) and \(\mathrm{pr}\) admit no linear
extension. The still point \(G_2(2)\) is transitive on the 56 weights,
on the 28 bitangents, on the 36 even theta characteristics and on the 63
nonzero vectors of \(\mathbb{F}_2^6\), with stabilizers of orders 216,
432, 336, 192; each even theta characteristic's stabilizer is a
\(PGL(2,7)\) whose derived group is one of the 36 subgroups of
Proposition 8.2, bijectively.} (14 checks.)

The fraction of weight pairs whose inner-product type is preserved is
\(100\%\) for \(\iota\) and for every Weyl element, \(93.0\%\) for
\(\mathrm{pr}\) and \(65.1\%\) for \(R\) (Proposition 4.1) --- the
observation from which §§3--4 grew.

\textbf{8.4 The fourteen outer classes.} The conjugacy classes of
\(O_8^+(2){:}3\) in the outer coset \(\Omega\Phi\) are classical data
{[}ATLAS; Bre, CTblLib 1.3.11{]}. The model re-derives them by
centralizer closures from the explicit permutation group of Proposition
8.3:

{\def\LTcaptype{none} % do not increment counter
\begin{longtable}[]{@{}llll@{}}
\toprule\noalign{}
order & classes & centralizer orders in \(\Omega\) & density \\
\midrule\noalign{}
\endhead
\bottomrule\noalign{}
\endlastfoot
24 & 2 & 8, 8 & \(1/8 + 1/8\) \\
21 & 1 & 7 & \(1/7\) \\
18 & 1 & 6 & \(1/6\) \\
12 & 4 & 32, 4, 48, 96 & \(5/16\) \\
9 & 1 & 18 & \(1/18\) \\
6 & 3 & 192, 24, 48 & \(13/192\) \\
3 & 2 & 12,096 (the turn), 216 & \(1/12{,}096 + 1/216\) \\
\end{longtable}
}

The densities sum to 1 exactly, so there is no fifteenth class
\textbf{(C)}. The turn's centralizer in \(\Omega\) is exactly
\(G_2(2)\); the element of order 18 has centralizer of order 6, and its
sixth power has centralizer of order \(1{,}944\) with centre of order 9.
Against the GAP Character Table Library's table of \(O_8^+(2).3\)
(headers only, no character values used) the outer class count 14 is
forced by the published class numbers (53 for \(\Omega\), 55 for
\(\Omega.3\), a fusion fixing 14 classes and fusing 13 triples); the
multiset of (order, centralizer order) of the published outer classes
equals ours tripled; all sixteen recorded power-map entries agree; the
sixth power of the eighteen is class \(3D\) of \(O_8^+(2)\), the \(E_8\)
and \(D_4\) Coxeter hearts are the fused triple \(3A/3B/3C\), and the
two classes of order 24 are told apart by their squares exactly as the
table tells them apart \textbf{(C)}.

\section{9. The bitangent bridge}\label{the-bitangent-bridge}

\textbf{9.1 \(PSL(2,7)\).} \(PSL(2,7) \cong GL(3,\mathbb{F}_2)\), order
168, classes \(1A, 2A, 3A, 4A, 7A, 7B\) of sizes
\(1, 21, 56, 42, 24, 24\), character table

{\def\LTcaptype{none} % do not increment counter
\begin{longtable}[]{@{}lllllll@{}}
\toprule\noalign{}
& \(1A\) & \(2A\) & \(3A\) & \(4A\) & \(7A\) & \(7B\) \\
\midrule\noalign{}
\endhead
\bottomrule\noalign{}
\endlastfoot
\(\chi_1\) & 1 & 1 & 1 & 1 & 1 & 1 \\
\(\chi_3\) & 3 & \(-1\) & 0 & 1 & \(\alpha\) & \(\bar\alpha\) \\
\(\bar\chi_3\) & 3 & \(-1\) & 0 & 1 & \(\bar\alpha\) & \(\alpha\) \\
\(\chi_6\) & 6 & 2 & 0 & 0 & \(-1\) & \(-1\) \\
\(\chi_7\) & 7 & \(-1\) & 1 & \(-1\) & 0 & 0 \\
\(\chi_8\) & 8 & 0 & \(-1\) & 0 & 1 & 1 \\
\end{longtable}
}

with \(\alpha = (-1 + \sqrt{-7})/2\) (re-derived by the Burnside--Dixon
class-algebra method and verified exactly over \(\mathbb{Q}(\sqrt{-7})\)
\textbf{(C)}). Since \(7 \nmid |W(E_6)| = 51{,}840\) and
\(168/27 \notin \mathbb{Z}\), \(PSL(2,7)\) embeds in no
\(W(E_6)\)-geometry and has no transitive action on 27 points. At field
level,
\(\mathrm{Gal}(\mathbb{Q}(\zeta_{21})/\mathbb{Q}) \cong C_6 \times C_2\)
has quadratic subfields \(\mathbb{Q}(\sqrt{-3})\),
\(\mathbb{Q}(\sqrt{-7})\), \(\mathbb{Q}(\sqrt{21})\);
\(\mathbb{Q}(\sqrt{-7})\) is the character field of \(PSL(2,7)\) and the
CM field of the Klein quartic {[}Elk99{]}, \(\mathbb{Q}(\sqrt{-3})\) the
Eisenstein field of \(E_6\); they are linearly disjoint.

\textbf{9.2 Theta characteristics.} Realise
\(PSL(2,7) = GL(3,\mathbb{F}_2)\) on
\(V \oplus V^* \cong \mathbb{F}_2^6\) with the symplectic form

\[\omega\big((v,f),(w,h)\big) = f(w) + h(v) .\]

The 64 quadratic refinements of \(\omega\) split by Arf invariant into
36 even and 28 odd; \(Sp_6(2)\) acts on them with the stabilizer of an
even form \(O_6^+(2) \cong S_8\) and of an odd form
\(O_6^-(2) \cong W(E_6)\) {[}Dol12, ch.~6{]}.

\textbf{Proposition 9.1 (C).} \emph{(i) \(PSL(2,7)\) fixes exactly one
even refinement; its orbits on the 36 even refinements have sizes
\(1, 7, 7, 21\), and it sits inside the stabilizer \(S_8\) of the fixed
form by its action on \(\mathbb{P}^1(\mathbb{F}_7)\). (ii) \(PSL(2,7)\)
is transitive on the 28 odd refinements with stabilizer
\(S_3 = N(\langle z_3 \rangle)\), the normaliser of a subgroup of order
3, and permutation character \(\chi_1 + 2\chi_6 + \chi_7 + \chi_8\).
(iii) Inside \(Sp_6(2) \cong W(E_7)/\{\pm 1\}\): \(W(E_6)\) is the
stabilizer of one odd refinement (one bitangent), \(PSL(2,7)\) is
transitive on all 28, and \(PSL(2,7) \cap \mathrm{Stab} = S_3\). (iv)
\(C_{Sp_6(2)}(PSL(2,7)) = 1\) and \(N_{Sp_6(2)}(PSL(2,7)) = PGL(2,7)\),
the outer automorphism realised as the symplectic duality
\(X \leftrightarrow X^*\) (by an exhaustive pass over all
\(1{,}451{,}520\) elements, and independently by a structural argument).
(v) The map \(z \mapsto N(\langle z \rangle)\) is exactly 2-to-1 from
the 56 elements of class \(3A\) onto the 28 bitangent stabilizers, and
the lift of \(PSL(2,7)\) to \(W(E_7)\) acts on the 56 weights with two
orbits of size 28, interchanged by \(\iota\), with weight stabilizer
\(S_3\); the two lifts to \(W(E_7) = 2 \times Sp_6(2)\) coincide because
\(PSL(2,7)\) is perfect.}

The \(S_3\) of (iii) is the stabilizer both of a bitangent in
\(W(E_6)\)'s geometry and of a point in \(PSL(2,7)\)'s action on 28
points --- the two subgroups are the two points of view on one
configuration. Two canonical restrictions of the 27-dimensional
representation of \(E_6\) follow from the two natural embeddings:
through \(E_6 \supset SU(3)^3\) with \(PSL(2,7)\) diagonal by
\(\chi_3\), \(27|_{PSL(2,7)} = 3\chi_1 \oplus 3\chi_8\); through
\(PSL(2,7) \subset G_2 = \mathrm{Aut}(\mathbb{O})\) with
\(27 = J_3(\mathbb{O})\), \(27| = 6\chi_1 \oplus 3\chi_7\) \textbf{(C)};
no restriction of the 27 contains a 16-dimensional constituent.

\textbf{9.3 The Clifford reading.} \(Sp_6(2) = W(E_7)/\{\pm1\}\) is the
3-qubit Clifford group modulo phases and Paulis; in that dictionary the
28 bitangents are the 28 odd Pauli sign-functions, \(W(E_6)\) the
Clifford stabilizer of one, and \(PSL(2,7)\) a Clifford subgroup
transitive on all 28 (exhibited computationally in shared coordinates,
35 checks \textbf{(C)}). \(R\) has no image in this quotient:
\(\iota R \iota = R^{-1}\) while \(\iota\) is central, and \(R\) does
not act on the 28 antipodal pairs (Corollary 5.2(iii)) --- the
group-theoretic face of Theorem A.

\textbf{9.4 The sextet and the Fano plane.}

\textbf{Proposition 9.2.} \emph{\(\chi_6\) is the deleted permutation
module on the 7 points of the Fano plane: realising
\(PSL(2,7) = GL(3,\mathbb{F}_2)\) on the 7 nonzero vectors of
\(\mathbb{F}_2^3\), generators \(A\) (order 2), \(B\) (order 3) with
\(AB\) of order 7 have traces \(2, 0, -1\) on the 6-dimensional sum-zero
subspace, which is \(\chi_6\)'s character at \(2A, 3A, 7A/7B\)
(cf.~{[}LNR07, §7{]}). The two orbits of size 7 in Proposition 9.1(i)
both carry permutation character \(1 \oplus \chi_6\) (fixed-point counts
\(7, 3, 1, 1, 0, 0\)): they are the points and the lines of one Fano
plane, dual under the outer automorphism, each realising \(\chi_6\) as
its deleted permutation module.} \textbf{(C)}

Consequently
\(\mathbb{C}\Lambda = 2(\chi_1 + 2\chi_6 + \chi_7 + \chi_8)\) contains
the Fano plane's geometry four times over and the three-dimensional
representations not at all; the octet \(\chi_8\) is the Steinberg
representation, and the 2-transitive action of \(PGL(2,7)\) on
\(\mathbb{P}^1(\mathbb{F}_7)\) appears as the stabilizer \(S_8\) of the
fixed even theta characteristic.

\section{10. Open problems}\label{open-problems}

\begin{enumerate}
\def\labelenumi{\arabic{enumi}.}
\tightlist
\item
  \textbf{Closed forms for \(\kappa\)} on the half-spins, on \(E_6\) and
  \(E_7\), and on \(A_n\omega_k\) for \(k \ge 3\) (Remark 3.6 suggests
  degree \(2k-2\) in \(n\)).
\item
  \textbf{The centralizer law.} Proposition 5.4 is a computation on 23
  boards. Is \(I_k = C_W(R^k)\) on every simply-laced minuscule board on
  which no power of \(R\) is a Weyl element, apart from the \(D_5\)
  half-spin? On the boards with a \(W\)-central reverser, Lemma 5.6
  gives \(I_k \subseteq C_W(R^{2k})\), with equality on the vector
  boards (Theorem 5.7); the inclusion is strict on \(A_5\omega_3\),
  \(A_7\omega_4\) and the \(D_6\) half-spins at the half-turn. Is there
  a description of \(I_k\) inside \(C_W(R^{2k})\) on those boards?
\item
  \textbf{The spinor boards.} Among \(B_3, \dots, B_8\) the shared
  grammar exceeds the centralizer only on \(B_4\) and, through
  \(W(D_4)\), on \(B_3\) (Proposition 6.7). Is \(B_4\) the only \(B_n\)
  spinor board with \(n \ge 4\) on which it does? The clock is free
  exactly when \(n\) is a power of 2, and \(B_{16}\) (65,536 weights) is
  the next free case; on \(B_{16}\) the correction's centralizer is
  again \(\{\pm 1\}\) (Proposition 6.11), and Theorem 6.14 shows that
  \(-(e_2e_3)(e_4e_5)\cdots\) commutes with \(R^{h/2}\) only at
  \(n = 4\). What remains is to show that this element is the Weyl
  element nearest to the half-turn for every \(n > 8\), and that
  \(C_W(R^{h/2}) = \{\pm 1\}\) for every \(n \ge 5\) (computed for
  \(n \le 8\) and \(n = 16\)).
\item
  \textbf{Non-simply-laced and non-minuscule boards}, promotion, and
  birational rowmotion {[}Oka21{]}: none is treated here.
\item
  \textbf{\(\langle W, R \rangle\)} on the boards where it could not be
  enumerated: it is \(A_{16}\) on the \(D_5\) half-spin and \(A_{56}\)
  on \(E_7\); ``alternating'' is an observation on five boards.
\item
  \textbf{The transports.} In which group below \(\mathrm{Sym}(2n)\), if
  any, are the outer transports of the odd \(D_n\) vector boards inner
  (Proposition 5.8)?
\end{enumerate}

\section{Acknowledgements}\label{acknowledgements}

The computations, the verification discipline (briefs hashed before
code, first runs kept) and much of the drafting were carried out with
Claude (Anthropic) as a computational assistant; the mathematics, the
identifications and the decisions of what to claim are the authors'. We
thank the authors of GAP and of the Character Table Library, against
whose published tables §8 was checked.

\section{Data and code availability}\label{data-and-code-availability}

Every statement marked \textbf{(C)} is established by a named Python 3
script (numpy; sympy in one place) listed in Appendix B. The scripts,
their hash-locked briefs, their logs including first runs, the witnesses
they write, and the two editions of this paper are kept under version
control at

\begin{quote}
\texttt{https://github.com/selinaserephina-star/Scar\_merkabit}
\end{quote}

The repository is private during review; access is granted to editors
and referees on request, and it will be made public on publication. Each
script re-checks its brief's SHA-256 as its first test and runs from a
clean checkout with no external data.

\section{Declarations}\label{declarations}

\emph{Funding.} This work received no external funding. \emph{Competing
interests.} The authors declare none. \emph{Author contributions.} S.S.
built the computational model and the verification framework, ran the
computations, and drafted the paper. I.B. built the \(PSL(2,7)\)
framework of §9, re-executed every computation independently, conducted
the literature search, and reviewed the text. Both authors approved the
final version.

\section{Appendix A. The literature
search}\label{appendix-a.-the-literature-search}

The claims of novelty in this paper are negative claims --- ``not
found'' --- and are worth exactly the search behind them. The question:
is it stated or proved anywhere that rowmotion on a minuscule poset,
carried to the weights, is or is not an element of the Weyl group; for
which posets it or its half-turn is; that a pair of weights keeps its
inner product under rowmotion according to a toggle parity; or what
fraction of pairs it keeps?

\emph{Keyword and citation-trail searches.} Toggle groups, periodicity,
homomesy, cyclic sieving, rowmotion with quadratic or bilinear forms mod
2, the \(E_7\) case and the Gosset graph; and the thirty documents
citing {[}RS13{]} on zbMATH Open, all read. Found: the Coxeter-motion
and birational literature {[}Oka21{]}, the toggle literature {[}SW12,
Str18{]}, the homomesy literature {[}PR15, RW15, DHPP23{]}; no match.

\emph{Corpus.} The arXiv API query \texttt{all:rowmotion} returns the
complete corpus of 57 documents (2014--2026). Every title and abstract
was read. Titles containing \emph{minuscule}: Hopkins 2016 (CDE property
for minuscule lattices), Hopkins 2019 (minuscule doppelgängers), Okada
{[}Oka21{]}, Pechenik {[}Pec22{]}. Titles containing \emph{Coxeter}:
Okada {[}Oka21{]}, Marczinzik--Thomas--Yıldırım {[}MTY24{]}. Titles
containing \emph{Weyl}, \emph{isometry}, \emph{orthogonal}, \emph{inner
product} or \emph{\(E_7\)}: none. The remaining titles concern Tamari
and Cambrian lattices, fences, root posets and rowvacuation,
interval-closed sets, products of chains, birational and
piecewise-linear lifts, Markov chains, rooted trees, trapezoids, plane
partitions, semidistributive and trim lattices, alternating sign
matrices, and echelonmotion. The zbMATH Open API returned nine documents
for the same term, none with these terms in title or keywords.

\emph{Primary texts, read for the question.} {[}RS13{]}: Theorem 1.3,
rowmotion and the toggle product \(t_{(i_1,\dots,i_n)}\) are conjugate
in the toggle group; Theorem 1.4, under Stembridge's bijection
\(J(P) \to W^J\) that product corresponds to a Coxeter element; no
statement that rowmotion itself is or is not a Coxeter element's action;
the chain case attributed to Stanley {[}Sta09{]}; inner products,
isometries and pairs of weights not mentioned. {[}Oka21{]}:
Coxeter-motion defined as a product of file toggles and shown conjugate
to rowmotion in the birational toggle group (Thm 4.3) --- two maps,
conjugate, not equal; the double-tailed diamond \(D_n\omega_1\) treated
birationally with separate formulas for the two tails, without any
statement about Weyl group membership. {[}Hop24{]}: records the
Rush--Shi conjugacy and the open conjectures on cyclic sieving and
doppelgängers; nothing on Weyl membership or pairs. {[}TW19{]}: trim
lattices; nothing on the weights.

\emph{Adjacent results, not the same statement.} {[}Oka21{]} and
{[}RS13{]} both relate rowmotion to a Coxeter-type action by conjugacy
and neither asks when the conjugating element can be taken trivial
(Theorems 3.2--3.3). {[}MTY24{]} study a homological ``Coxeter matrix''
built from the incidence algebra's Cartan matrix alongside rowmotion and
prove a different commutation identity. {[}Pan16{]} indexes lower ideals
of weight posets by Weyl group elements, not rowmotion by one.

\emph{What the search does not reach.} MathSciNet and the full text of
zbMATH's reviews; theses and work in progress. A reader who knows
Theorems A--G as corollaries of the Rush--Shi conjugacy or of the toggle
description of rowmotion is asked for the reference.

\section{Appendix B. Verification}\label{appendix-b.-verification}

Every statement marked (C) is established by a Python 3 script (numpy;
sympy in one place) whose brief was fixed and hashed before the script
was written and which re-checks that hash as its first test. Each script
was executed by both authors from clean checkouts on independent
machines with identical tallies. The scripts, with their check counts
and external inputs:

All scripts are at the root of the repository except
\texttt{verify\_tbr\_bridge.py} and \texttt{verify\_tsc\_scarcat.py}
(under \texttt{package/}) and \texttt{fano\_sextet\_six\_invariants.py}
(in the received-packages folder). ``Inputs'' are files written by the
scripts of earlier rows; all are regenerable.

{\def\LTcaptype{none} % do not increment counter
\begin{longtable}[]{@{}
  >{\raggedright\arraybackslash}p{(\linewidth - 6\tabcolsep) * \real{0.2500}}
  >{\raggedright\arraybackslash}p{(\linewidth - 6\tabcolsep) * \real{0.2500}}
  >{\raggedright\arraybackslash}p{(\linewidth - 6\tabcolsep) * \real{0.2500}}
  >{\raggedright\arraybackslash}p{(\linewidth - 6\tabcolsep) * \real{0.2500}}@{}}
\toprule\noalign{}
\begin{minipage}[b]{\linewidth}\raggedright
statement
\end{minipage} & \begin{minipage}[b]{\linewidth}\raggedright
script
\end{minipage} & \begin{minipage}[b]{\linewidth}\raggedright
checks
\end{minipage} & \begin{minipage}[b]{\linewidth}\raggedright
inputs
\end{minipage} \\
\midrule\noalign{}
\endhead
\bottomrule\noalign{}
\endlastfoot
Prop. 8.1 & \texttt{verify\_stone\_v\_we8\_roof.py} & 44 & cache of
\texttt{verify\_stone\_u\_2cover.py} \\
Prop. 8.2 & \texttt{verify\_stone\_z\_stillpoint.py} & 19 + 1 inverted &
caches of the \texttt{u}, \texttt{v}, \texttt{w}, \texttt{x} scripts \\
Prop. 8.3 & \texttt{verify\_stone\_aa\_roofclock.py} & 16 & the same,
plus the cache of \texttt{verify\_stone\_z\_stillpoint.py} \\
Prop. 8.4 & \texttt{verify\_stone\_ab\_merkabit.py} & 14 &
\texttt{scar56\_data.json}, \texttt{G\_rows.npy} \\
Prop. 4.1 (E₇), §4.1 & \texttt{verify\_stone\_ac\_parity\_rule.py} & 9 +
2 inverted & \texttt{scar56\_data.json} \\
§8.4 (14 classes) & \texttt{verify\_stone\_ai..am\_*.py} & 15, 17, 14,
15+2, 15+1 & caches of the \texttt{aa} and \texttt{z} scripts \\
§8.4 (ATLAS names) & \texttt{verify\_stone\_an\_atlas\_names.py} & 17 &
+ two CTblLib header files \\
Prop. 4.4 & \texttt{verify\_stone\_ao\_mirror\_and\_gates.py} & 11 + 1
inverted & \texttt{scar56\_data.json} \\
Prop. 4.2 & \texttt{verify\_stone\_ap\_clock\_centralizer.py} & 22 + 2 &
\texttt{scar56\_data.json} \\
Table 1 & \texttt{verify\_stone\_aq\_rush\_shi\_defect.py} & 8 + 2
inverted & none (Cartan matrices) \\
Thms 3.2--3.5 checks &
\texttt{verify\_stone\_ar\_nonlinearity\_theorems.py} & 10 & the board
machine of the \texttt{aq} script \\
Thm 4.3 & \texttt{verify\_stone\_as\_shared\_grammar.py} & 14 &
witnesses of the \texttt{ap} script, \texttt{scar56\_data.json} \\
Prop. 5.4 & \texttt{verify\_stone\_at\_grammar\_family.py} & 7 + 4 & the
board machine of the \texttt{aq} script \\
Thm 5.1, Prop. 4.5 & \texttt{verify\_stone\_au\_two\_reversals.py} & 10
& witnesses of the \texttt{aq}, \texttt{ap} scripts \\
§6.1 & \texttt{verify\_stone\_av\_d5\_anomaly.py} & 9 + 4 & none \\
Prop. 7.1 & \texttt{verify\_stone\_aw\_homomesy.py} & 6 + 1 & the board
machine of the \texttt{aq} script \\
Props. 4.6, 5.3 & \texttt{verify\_stone\_ax\_c\_equals\_t.py} & 21 + 1 &
witnesses of the \texttt{aq}, \texttt{at}, \texttt{ap}, \texttt{as},
\texttt{au} scripts \\
Thm 6.2, Lemma 6.1 & \texttt{verify\_stone\_ay\_d5\_by\_hand.py} & 20 +
5 & witnesses of the \texttt{av}, \texttt{aq}, \texttt{at} scripts \\
Thm 5.7, Prop. 5.8 & \texttt{verify\_stone\_az\_doubled\_clock.py} & 11
+ 2 & witnesses of the \texttt{av}, \texttt{aq}, \texttt{at} scripts \\
Prop. 6.6 & \texttt{verify\_stone\_ba\_clifford\_hierarchy.py} & 7 &
witnesses of the \texttt{aq}, \texttt{at}, \texttt{av} scripts \\
Thm 6.2 (rank-shift) &
\texttt{verify\_stone\_bb\_halfturn\_fullheight.py} & 6 & none \\
Prop. 6.7 & \texttt{verify\_stone\_bc\_b8\_halfturn.py} & 14 + 4
inverted & none \\
Lemma 6.9, Prop. 6.10 &
\texttt{verify\_stone\_bd\_halfturn\_criterion.py} & 7 + 3 inverted &
none \\
Prop. 6.11 & \texttt{verify\_stone\_be\_b16\_orbits.py} & 9 + 2 inverted
& none \\
Lemma 6.12, Prop. 6.13, Thm 6.14 &
\texttt{verify\_stone\_bf\_why\_four.py} & 8 & none \\
Prop. 9.1 & \texttt{verify\_tbr\_bridge.py},
\texttt{verify\_stonef\_b\_dq1.py} & 20, 11 & none \\
§9.1 table & \texttt{verify\_tsc\_scarcat.py} & 48 & none \\
§9.3 & \texttt{verify\_stoneq\_clifford.py} & 35 & none \\
Prop. 9.2 & \texttt{fano\_sextet\_six\_invariants.py},
\texttt{verify\_fano\_bitangent\_identity.py} & re-run, no tally &
none \\
\end{longtable}
}

``Inverted'' counts record expectations registered before the
computation that the computation refuted; each is described where it is
informative (Remarks 3.7, 5.5, 5.9; Propositions 8.2, 8.4). The scripts,
the hash-locked briefs and the logs are in the repository named under
\emph{Data and code availability}.

\section{References}\label{references}

\begin{itemize}
\tightlist
\item
  {[}ATLAS{]} Conway, J. H.; Curtis, R. T.; Norton, S. P.; Parker, R.
  A.; Wilson, R. A. \emph{ATLAS of Finite Groups.} Clarendon Press,
  Oxford University Press, 1985.
\item
  {[}Bou02{]} Bourbaki, N. \emph{Lie Groups and Lie Algebras, Chapters
  4--6.} Elements of Mathematics, Springer, Berlin, 2002.
\item
  {[}Bre25{]} Breuer, T. \emph{The GAP Character Table Library}, Version
  1.3.11 (GAP package), Lehrstuhl für Algebra und Zahlentheorie, RWTH
  Aachen, 2025.
  \texttt{https://www.math.rwth-aachen.de/\textasciitilde{}Thomas.Breuer/ctbllib}
\item
  {[}CF95{]} Cameron, P. J.; Fon-der-Flaass, D. G. Orbits of antichains
  revisited. \emph{European J. Combin.} 16 (1995), no. 6, 545--554.
\item
  {[}Con71{]} Conway, J. H. Three lectures on exceptional groups. In
  \emph{Finite Simple Groups} (M. B. Powell, G. Higman, eds.), Academic
  Press, London--New York, 1971, 215--247.
\item
  {[}CS99{]} Conway, J. H.; Sloane, N. J. A. \emph{Sphere Packings,
  Lattices and Groups}, 3rd ed.~Grundlehren Math. Wiss. 290, Springer,
  New York, 1999.
\item
  {[}DHPP23{]} Defant, C.; Hopkins, S.; Poznanović, S.; Propp, J.
  Homomesy via toggleability statistics. \emph{Combinatorial Theory} 3
  (2023), no. 2, \#14; arXiv:2108.13227.
\item
  {[}Dol12{]} Dolgachev, I. \emph{Classical Algebraic Geometry: A Modern
  View.} Cambridge University Press, 2012.
\item
  {[}Elk99{]} Elkies, N. D. The Klein quartic in number theory. In
  \emph{The Eightfold Way} (S. Levy, ed.), MSRI Publ. 35, Cambridge
  University Press, 1999, 51--101.
\item
  {[}Gre13{]} Green, R. M. \emph{Combinatorics of Minuscule
  Representations.} Cambridge Tracts in Mathematics 199, Cambridge
  University Press, 2013.
\item
  {[}GX26{]} Green, R. M.; Xu, T. Branching rules of minuscule
  representations via a new partial order. \emph{Combinatorial Theory} 6
  (2026), no. 1, \#2; arXiv:2402.06732.
\item
  {[}Hop24{]} Hopkins, S. Order polynomial product formulas and poset
  dynamics. In \emph{Open Problems in Algebraic Combinatorics}, Proc.
  Sympos. Pure Math. 110, Amer. Math. Soc., Providence, RI, 2024,
  135--157; arXiv:2006.01568.
\item
  {[}Kle87{]} Kleidman, P. B. The maximal subgroups of the finite
  8-dimensional orthogonal groups \(P\Omega_8^+(q)\) and of their
  automorphism groups. \emph{J. Algebra} 110 (1987), no. 1, 173--242.
\item
  {[}LNR07{]} Luhn, C.; Nasri, S.; Ramond, P. Simple finite non-abelian
  flavor groups. \emph{J. Math. Phys.} 48 (2007), no. 12, 123519;
  arXiv:0709.1447.
\item
  {[}MTY24{]} Marczinzik, R.; Thomas, H.; Yıldırım, E. On the
  interaction of the Coxeter transformation and the rowmotion bijection.
  \emph{J. Comb. Algebra} 8 (2024), no. 3/4, 359--374; arXiv:2201.04446.
\item
  {[}Oka21{]} Okada, S. Birational rowmotion and Coxeter-motion on
  minuscule posets. \emph{Electron. J. Combin.} 28 (2021), no. 1,
  \#P1.17; arXiv:2004.05364.
\item
  {[}Pan16{]} Panyushev, D. I. Weight posets associated with gradings of
  simple Lie algebras, Weyl groups, and arrangements of hyperplanes.
  \emph{J. Algebraic Combin.} 44 (2016), no. 2, 325--344;
  arXiv:1412.0987.
\item
  {[}Pec22{]} Pechenik, O. Minuscule analogues of the plane partition
  periodicity conjecture of Cameron and Fon-Der-Flaass.
  \emph{Combinatorial Theory} 2 (2022), no. 1, \#15; arXiv:2107.02679.
\item
  {[}Pro84{]} Proctor, R. A. Bruhat lattices, plane partition generating
  functions, and minuscule representations. \emph{European J. Combin.} 5
  (1984), no. 4, 331--350.
\item
  {[}PR15{]} Propp, J.; Roby, T. Homomesy in products of two chains.
  \emph{Electron. J. Combin.} 22 (2015), no. 3, \#P3.4.
\item
  {[}RSW04{]} Reiner, V.; Stanton, D.; White, D. The cyclic sieving
  phenomenon. \emph{J. Combin. Theory Ser. A} 108 (2004), no. 1, 17--50.
\item
  {[}RS13{]} Rush, D. B.; Shi, X. On orbits of order ideals of minuscule
  posets. \emph{J. Algebraic Combin.} 37 (2013), no. 3, 545--569;
  arXiv:1108.5245.
\item
  {[}RW15{]} Rush, D. B.; Wang, K. On orbits of order ideals of
  minuscule posets II: Homomesy. arXiv:1509.08047 (2015).
\item
  {[}Sta09{]} Stanley, R. P. Promotion and evacuation. \emph{Electron.
  J. Combin.} 16 (2009), no. 2, \#R9.
\item
  {[}Ste94{]} Stembridge, J. R. On minuscule representations, plane
  partitions and involutions in complex Lie groups. \emph{Duke Math. J.}
  73 (1994), no. 2, 469--490.
\item
  {[}Ste01{]} Stembridge, J. R. Minuscule elements of Weyl groups.
  \emph{J. Algebra} 235 (2001), no. 2, 722--743.
\item
  {[}Str18{]} Striker, J. Rowmotion and generalized toggle groups.
  \emph{Discrete Math. Theor. Comput. Sci.} 20 (2018), no. 1, Paper
  No.~17; arXiv:1601.03710.
\item
  {[}SW12{]} Striker, J.; Williams, N. Promotion and rowmotion.
  \emph{European J. Combin.} 33 (2012), no. 8, 1919--1942.
\item
  {[}TW19{]} Thomas, H.; Williams, N. Rowmotion in slow motion.
  \emph{Proc. London Math. Soc. (3)} 119 (2019), no. 5, 1149--1178;
  arXiv:1712.10123.
\item
  {[}Wil09{]} Wilson, R. A. \emph{The Finite Simple Groups.} Graduate
  Texts in Mathematics 251, Springer, London, 2009.
\end{itemize}

\begin{center}\rule{0.5\linewidth}{0.5pt}\end{center}

Selina Stenberg --- Independent researcher ---
\texttt{selina@exoreaction.com}

Ilya Balashov --- Independent researcher --- ORCID 0009-0001-1954-8872

\end{document}
